% SPDX-License-Identifier: CC-BY-4.0
% =====================================================================
% Section 7 of D. Joyner and Claude, "A hitch-hiker's guide to the
% metaplectic trace formula for SL(2)", revised for the 2-fold cover
% (n = 2), November 2026.
%
% Stage 6 of ../revision_log.md. Sources: ../original-1995/cctrace7.tex
% (= ../book-1997/book7h.tex) and ../book-1997/book3b.tex. Referee items
% II.18-II.23, C.1, C.3, C.4 and the Section 7 part of I.1 and II.22 are
% treated here; see the "Stage 6 notes" in ../revision_log.md, which also
% map the 1995 numbers to the numbers here.
%
% This file is \input by the master; it compiles on its own through
% sl2bar-trace-2026.tex. Explicit formulas are checked by checks/sec7_n2_check.py.
% Comments beginning "% [Claude]" are editorial notes by Claude.
% =====================================================================

\providecommand{\GA}{\overline{G(\A)}}
\providecommand{\GS}{\overline{G(\Q_S)}}
\providecommand{\AS}{\overline{A(\Q_S)}}
\providecommand{\KS}{\overline{K_S}}
\providecommand{\Gb}{\overline{SL(2,F)}}

\section{The trace formula in stable(r) form}\label{sec:stable}

The invariant trace formula of Theorem~\ref{thm:itf} is an identity
between distributions that are invariant under conjugation by the cover
itself. For $SL(2)$, Labesse and Langlands~\cite{LL} went one step
further: they regrouped the terms according to conjugacy under $GL(2)$,
which for $SL(2)$ is stable conjugacy, and identified the part of the
formula that is not stably invariant with contributions of the elliptic
tori. This section carries out the first half of that program for the
double cover: it determines how $GL(2)$ acts on the cover, what stable
conjugacy is, and how each term of the trace formula behaves. We call the
result the \emph{stable(r) form}: the terms are organized by stable
classes, and the elliptic regular terms are written as stable orbital
integrals, but nothing is transferred to another group.

Three properties of a distribution on the cover occur below and should be
kept apart. \emph{Equivariance}: it is multiplied by a character of $\det
h$ when $f$ is replaced by $f^h=f\circ\mathrm{Ad}(h)$, $h\in GL(2)$; this
is about ordinary conjugation, and it is what
Section~\ref{sec:hyperbolic-gl2} proves for the weighted terms.
\emph{Stability}: it is invariant under the twisted action $f\mapsto
f^{[h]}$ of Section~\ref{sec:stable-conj}; this is proved for the regular
elliptic terms, the characters of the principal series and the theta
functions. \emph{Endoscopic transfer}: it is identified with a stable
distribution on $SO(3)$ applied to a transfer of $f$; this is Li's
theorem, and nothing of the kind is proved here.

For the double cover of $Sp(2n)$ the full theory exists and is due to
Li: stable conjugacy and the transfer of orbital
integrals~\cite{Li11}, following Adams~\cite{Ad}, Renard~\cite{Re} and,
for the group here over a $p$-adic field, Schultz~\cite{Sch} (see
also~\cite{How}); the stabilization of the elliptic
terms~\cite{Li15}; and the stabilization of the whole invariant trace
formula~\cite{Li21}. The elliptic endoscopic groups of the double cover of
$Sp(2n)$ are the groups $SO(2n'+1)\times SO(2n''+1)$, $n'+n''=n$. For
$n=1$ these are two copies of $SO(3)=PGL(2)$, attached to the pairs
$(1,0)$ and $(0,1)$; the tori of $SL(2)$, which are the endoscopic groups
of $SL(2)$ in~\cite{LL}, do not occur. We do not reprove any of this.
What can be done by hand for the group of this paper, and is done below,
is the following.

\paragraph{Results.}
\begin{itemize}
\item $GL(2,F)$ acts on $\Gb$ by automorphisms $\mathrm{Ad}(h)$, with an
  explicit formula (Lemma~\ref{lem:Ad}). An element $h$ that commutes
  with a semisimple $\gamma$ multiplies the lifts of $\gamma$ by the sign
  $(\det h,\det(1+\gamma))_v$ (Lemma~\ref{lem:eta}). So conjugation by
  $GL(2)$ can carry a lift of $\gamma$ to the other lift, and it is not
  stable conjugacy.
\item Stable conjugacy on the cover is $GL(2)$-conjugacy \emph{twisted}
  by that sign: $h\star x=(\det h,\det(1+x))_v\,\mathrm{Ad}(h)x$. Its
  classes are the level sets, over the stable classes of $SL(2,F)$, of an
  explicit genuine function $\Delta_\psi$ with values in $\mu_4$
  (Lemma~\ref{lem:Delta}). This function is $|\det(1+x)|^{1/2}$ times the
  character of the difference $r^+_\psi-r^-_\psi$ of the two halves of
  the Weil representation (Proposition~\ref{prop:weil-char}), so the
  definition agrees with that of Adams and Li.
\item For a regular elliptic $\gamma\in G(\Q)$ with splitting field $E$,
  the terms of the trace formula attached to the $GL(2,k)$-conjugates of
  $\gamma$ add up to $\frac12$ of a stable orbital integral plus $\frac12$
  of a $\kappa$-orbital integral, as for $SL(2)$. For the cover the
  $\kappa$-orbital integral at $\gamma$ is again a stable orbital
  integral: that of the translate $f_-$ of $f$ by the central element
  $-1$, at $-\gamma$ (Lemma~\ref{lem:kappa-orbital}). Hence the elliptic
  regular part of the trace formula is $S_{\mathrm{ell}}(f)+
  S_{\mathrm{ell}}(f_-)$ with one stable distribution $S_{\mathrm{ell}}$
  (Theorem~\ref{thm:elliptic-stable}). This is the two-term pattern, with
  coefficients $\frac12$ and $\frac12$, of the elliptic part of Li's
  stable trace formula for $n=1$ with its two endoscopic data; the terms
  are not matched with the transferred distributions on $SO(3)$.
\item The remaining geometric terms transform under $GL(2,k_S)$ by
  explicit characters (Proposition~\ref{prop:IA-gl2}), and the
  unipotent coefficients $a^G(S,u_\xi)$, left open in
  Section~\ref{sec:geometric}, are computed:
  $a^G(S,u_\xi)=\frac12\bigl(c_S+\sum_{\kappa\ne1}\kappa_S(\xi)L^S(1,\kappa)
  \bigr)$, the sum over the quadratic Hecke characters unramified outside
  $S$ (Proposition~\ref{prop:unipotent}).
\item On the spectral side, $GL(2,k)$ permutes the pieces of the
  spectrum, and $V_b\mapsto V_{ab}$. The theta functions attached to
  $\psi_{-b}$ span a subspace of the discrete spectrum whose trace is
  $\frac12\Theta(f)+\frac12\Theta(f_-)$ with $\Theta$ a stable character
  (Proposition~\ref{prop:theta-packet}): the simplest instance of the
  same pattern on the spectral side.
\end{itemize}
The hyperbolic terms, the unipotent terms and the cuspidal spectrum are
\emph{not} stabilized here; Section~\ref{sec:stable-form} says what is
missing, and Remark~\ref{rem:sec7-status} what rests on citation.

\paragraph{Conventions.} $F=k_v$ is a local field of characteristic $0$,
as in Section~\ref{sec:local}; we write $G_1=GL(2)$, $Z$ for its center,
and $\delta_a=\diag(a,1)$. For $h\in G_1(F)$, $\mathrm{Ad}(h)$ is the
automorphism of $\Gb$ of Lemma~\ref{lem:gl2}(c), and for a function $f$
on the cover
\[
f^h(x)=f(\mathrm{Ad}(h)x),
\]
which for $h\in SL(2,F)$ is the function $f^y$ of
Section~\ref{sec:invariant}, $y$ a lift of $h$. For a central element
$z$ of a cover, $f_z(x)=f(zx)$, as in Section~\ref{sec:coarse}. The
element $-1\in G(\Q)$ is identified with its image $s(-1)=(-1,1)$ in
$\GA$, and
\[
f_-(x)=f_{-1}(x)=f\bigl(s(-1)\,x\bigr),\qquad f\in\mathcal H(\GA).
\]
Since $s(-1)^2=1$, $(f_-)_-=f$. For a quadratic extension $E=F(\sqrt\theta)$
let $\kappa_E=(\cdot,\theta)_v$ be the quadratic character of $F^\times$
with kernel $N E^\times=N_{E/F}(E^\times)$; for $E=F\times F$ put
$N E^\times=F^\times$ and $\kappa_E=1$.

\subsection{Conjugation by \texorpdfstring{$GL(2)$}{GL(2)}}\label{sec:Ad}

Conjugation by $h\in G_1(F)$ is an automorphism of $SL(2,F)$, and it
lifts to the cover in exactly one way (for the metaplectic cover of
$Sp(V)$ and the similitude group in general, see~\cite[\S3]{Th13}). This subsection makes the lift
explicit and computes it on the subgroups that matter: $N$, $A$, the
center, the maximal compact subgroup, the tori, and globally on $G(\Q)$.

\begin{lemma}\label{lem:Ad}
Let $h\in G_1(F)$.
\begin{enumerate}[label=(\alph*)]
\item $\mathrm{Ad}(h)$ is the only automorphism of the group $\Gb$ that
  induces $g\mapsto hgh^{-1}$ on $SL(2,F)$. It is a homeomorphism for
  the topology of Lemma~\ref{lem:kubota}(c), and it fixes $\mu_2$
  pointwise.
\item $\mathrm{Ad}(h)$ depends only on the image of $h$ in $PGL(2,F)$.
  For $h\in SL(2,F)$ it is conjugation by either lift of $h$. So
  $\mathrm{Ad}(h)$ is inner when $\det h\in F^{\times2}$, and the group
  $F^\times/F^{\times2}$ acts through the $\mathrm{Ad}(\delta_a)$ modulo
  inner automorphisms.
\item For $a\in F^\times$ and $g=\begin{pmatrix}\alpha&b\\c&d
  \end{pmatrix}\in SL(2,F)$,
  \begin{equation}\label{eq:Ad-delta}
  \mathrm{Ad}(\delta_a)(g,\zeta)=\bigl(\delta_ag\delta_a^{-1},\ v_a(g)\,
  \zeta\bigr),\qquad v_a(g)=\begin{cases}1,&c\ne0,\\(a,d)_v,&c=0.
  \end{cases}
  \end{equation}
\item In particular
  \begin{align*}
  \mathrm{Ad}(\delta_a)(n(u),1)&=(n(au),1),&
  \mathrm{Ad}(\delta_a)(a(t),\zeta)&=(a(t),(a,t)_v\zeta),\\
  \mathrm{Ad}(\delta_a)(-1,\zeta)&=(-1,(a,-1)_v\zeta),&
  \mathrm{Ad}(\delta_a)(w,1)&=(a(a),1)(w,1).
  \end{align*}
\item Let $v$ be finite and non-dyadic. If $h\in Z(F)GL(2,\OO)$, then
  $\mathrm{Ad}(h)K^*=K^*$ and $e_v^h=e_v$ for the unit element $e_v$
  of~\eqref{eq:unit}.
\item (Global.) For $h=(h_v)\in G_1(\A_k)$ the automorphisms
  $\mathrm{Ad}(h_v)$ define an automorphism $\mathrm{Ad}(h)$ of $\GA$,
  and $f^h\in\mathcal H(\GA)$ for $f\in\mathcal H(\GA)$. If
  $h\in G_1(k)$, then $\mathrm{Ad}(h)$ preserves the image of $G(\Q)$:
  $\mathrm{Ad}(h)s(\gamma)=s(h\gamma h^{-1})$ for $\gamma\in G(\Q)$, where
  $s(\gamma)=(\gamma,\kappa_\A(\gamma))$ is the splitting of
  Lemma~\ref{lem:global}(a).
\end{enumerate}
\end{lemma}

\begin{proof}
(a) Two automorphisms that induce the same automorphism of $SL(2,F)$
differ by a homomorphism $SL(2,F)\to\mu_2$, which is trivial because
$SL(2,F)$ is perfect. $\mathrm{Ad}(h)$ fixes $\mu_2$ because $\mu_2$ is
central in the cover of $GL(2,F)$. For continuity it suffices to look
near the identity. Let $U$ be as in Lemma~\ref{lem:kubota}(b), and let
$V\subset U$ be a neighborhood of $1$ with $hVh^{-1}\subset U$. For
$g\in V$ write $\mathrm{Ad}(h)(g,1)=(hgh^{-1},\varepsilon(g))$. Since
$\beta\equiv1$ on $U\times U$, $\varepsilon(g_1g_2)=\varepsilon(g_1)
\varepsilon(g_2)$ whenever $g_1,g_2,g_1g_2\in V$; and every $g$ in a
smaller neighborhood is a square of an element of $V$, as in the proof of
Lemma~\ref{lem:kubota}(c). So $\varepsilon=1$ near $1$, and
$\mathrm{Ad}(h)$ maps the basic neighborhoods $\{(g,1)\mid g\in V'\}$ to
basic neighborhoods.

(b) A scalar matrix induces the identity on $SL(2,F)$, so by~(a) it acts
trivially on the cover; and conjugation by a lift of $h\in SL(2,F)$ is an
automorphism over $g\mapsto hgh^{-1}$. If $\det h=c^2$, then $c^{-1}h\in
SL(2,F)$.

(c) By definition $\mathrm{Ad}(\delta_a)(g,\zeta)=(\delta_a,1)(g,\zeta)
(\delta_a,1)^{-1}$ in $\overline{GL(2,F)}$, which is
\[
\bigl(\delta_ag\delta_a^{-1},\ \zeta\,\sigma(\delta_a,g)\,\sigma(\delta_ag,
\delta_a^{-1})\,\sigma(\delta_a,\delta_a^{-1})\bigr)
\]
with the cocycle~\eqref{eq:sigma-gl2}. Here $x(\delta_a)=x(\delta_a^{-1})=1$ and
$x(\delta_ag)=x(g)$, so $\sigma(\delta_a,g)=(x(g),a)_v$ and
$\sigma(\delta_a,\delta_a^{-1})=1$. If $c\ne0$, then $x(\delta_ag
\delta_a^{-1})=c/a$ and $\sigma(\delta_ag,\delta_a^{-1})=(a^{-1},c/a^2)_v
=(a,c)_v$, so the product is $1$. If $c=0$, then
$x(\delta_ag\delta_a^{-1})=d=x(g)$ and $\sigma(\delta_ag,\delta_a^{-1})=
(1,d/a)_v=1$, so the product is $(d,a)_v$. (The script also checks
directly that~\eqref{eq:Ad-delta} is a homomorphism for $\beta$, which
by~(a) proves~(c) without reference to the cover of $GL(2)$.)

(d) The first three are~\eqref{eq:Ad-delta} with $c=0$ and $d=1$,
$t^{-1}$, $-1$. For the last, $\delta_aw\delta_a^{-1}=a(a)w$ has lower
left entry $a^{-1}\ne0$, so $\mathrm{Ad}(\delta_a)(w,1)=(a(a)w,1)$, and
$\beta(a(a),w)=(1,a^{-1})_v=1$.

(e) By~(b) we may take $h\in GL(2,\OO)$, so $hKh^{-1}=K$. Then
$k\mapsto\mathrm{Ad}(h)^{-1}(hkh^{-1},\kappa(hkh^{-1}))$ is a
homomorphic section of $\overline K\to K$. It differs from
$k\mapsto(k,\kappa(k))$ by a homomorphism $K=SL(2,\OO)\to\mu_2$, and there
is none except the trivial one: $K$ is generated by elementary matrices,
which are squares.

(f) By~(e), $\mathrm{Ad}(h_v)$ preserves $K_v^*$ for almost all $v$, so
the product acts on the restricted product of the local covers; it fixes
every $\mu_2$, hence the subgroup $Z_0$ of Section~\ref{sec:global}, and
descends to $\GA$. For $f=\bigotimes f_v$ we get
$f^h=\bigotimes f_v^{h_v}$ with $f_v^{h_v}=e_v$ for almost all $v$. If
$h\in G_1(k)$, then $\gamma\mapsto\mathrm{Ad}(h)^{-1}s(h\gamma h^{-1})$ is a
homomorphic section of $\overline{G(\Q)}\to G(\Q)$; it equals $s$ because
$SL(2,k)$ is perfect.
\end{proof}

Part~(d) already shows that $\mathrm{Ad}(h)$ is not inner in general: it
acts on the abelian group $\overline{A(F)}$ by the character
$a(t)\mapsto(a,t)_v$, and it exchanges the two lifts of $-1$ when
$(\det h,-1)_v=-1$. The next lemma gives the general statement. For a
semisimple $\gamma\in SL(2,F)$, $\gamma\ne-1$, put
\begin{equation}\label{eq:eta}
D^+(\gamma)=\det(1+\gamma)=2+\operatorname{tr}\gamma\ \in F^\times,\qquad
\eta_\gamma(a)=\bigl(a,D^+(\gamma)\bigr)_v\quad(a\in F^\times).
\end{equation}
The definition of $D^+$ and $\eta_\gamma$ makes sense for every $\gamma$
with $\operatorname{tr}\gamma\ne-2$, and depends only on the trace.

\begin{lemma}\label{lem:eta}
Let $\gamma\in SL(2,F)$ be semisimple, $\gamma\ne-1$, and let $h\in
G_1(F)$ commute with $\gamma$. Then
\[
\mathrm{Ad}(h)(\gamma,\zeta)=\bigl(\gamma,\ \eta_\gamma(\det h)\,\zeta
\bigr).
\]
For $\gamma=-1$ and every $h\in G_1(F)$,
$\mathrm{Ad}(h)(-1,\zeta)=(-1,(\det h,-1)_v\zeta)$.
\end{lemma}

\begin{proof}
The second statement is Lemma~\ref{lem:Ad}(b),(d). For $\gamma=1$ there is
nothing to prove ($D^+=4$). Let $\gamma$ be regular. Its centralizer in
$M_2(F)$ is the commutative algebra $E=F[\gamma]$, and $h\in E^\times$.
Put $z=1+\gamma\in E$. Then $\det z=D^+(\gamma)\ne0$, and $z\gamma^{-1}=
1+\gamma^{-1}$ is the conjugate $\bar z$ of $z$ in $E$, so
$z\bar z=\det z$ and
\[
\gamma\cdot\bigl(D^+(\gamma)\cdot1\bigr)=z^2\qquad\text{in }E^\times
\subset G_1(F).
\]
For commuting $x,y\in G_1(F)$ let $[x,y]\in\mu_2$ be the commutator of
any two lifts to $\overline{GL(2,F)}$. On the commutative group
$E^\times$ it is bimultiplicative, so $[h,z^2]=[h,z]^2=1$ and therefore
$[h,\gamma]=[h,D^+(\gamma)\cdot1]$. For a scalar $\lambda$ and any $h$,
\eqref{eq:sigma-gl2} gives $\sigma(h,\lambda)=(\lambda,x(h)/\det h)_v$ and
$\sigma(\lambda,h)=(x(h),\lambda)_v$, so $[h,\lambda]=(\lambda,\det h)_v$.
Hence $\mathrm{Ad}(h)(\gamma,\zeta)=(\gamma,[h,\gamma]\zeta)$ with
$[h,\gamma]=(D^+(\gamma),\det h)_v$.
\end{proof}

By~\eqref{eq:gl2-commutator} the sign is also $(x(h\gamma)/x(h),\det
h)_v$; the lemma says that this equals $(\det h,2+\operatorname{tr}
\gamma)_v$, which the script confirms on random elements of the tori
$T_\theta(\Q)$. If $\gamma$ is regular with centralizer $E^\times$ in
$G_1(F)$, then $\det h$ runs through $N E^\times$, and $D^+(\gamma)=
N_{E/F}(1+\gamma)$ is itself a norm. So the lifts of $\gamma$ are
exchanged by the centralizer of $\gamma$ in $G_1(F)$ exactly when
$\eta_\gamma$ is not trivial on $N E^\times$.

\begin{example}\label{ex:eta}
(a) \emph{Split torus.} For $\gamma=a(t)$, $t\ne\pm1$:
$D^+=(1+t)^2/t$, so $\eta_\gamma=(\cdot,t)_v$ and $N E^\times=F^\times$.
The two lifts of $a(t)$ are conjugate under $G_1(F)$ if and only if
$t\notin F^{\times2}$.

(b) \emph{$F=\R$, $\gamma$ elliptic.} $N E^\times=\R_{>0}$ and
$D^+=2+2\cos\theta>0$, so $\eta_\gamma=1$. In $Mp(2,\R)$ the automorphism
$\mathrm{Ad}(\delta_{-1})$ sends the element $\tilde k_\theta$ of
Lemma~\ref{lem:explicit-r}(b) to $\tilde k_{-\theta}$; on the center
$\{\tilde k_{j\pi}\}\cong\Z/4$ it is inversion, in agreement with
$(-1,-1)_\R=-1$.

(c) \emph{$F$ non-dyadic, $E/F$ unramified, $\gamma\in K$.} The norms are
the elements of even valuation, and $\eta_\gamma$ is trivial on them if
$D^+(\gamma)$ is a unit. But $\eta_\gamma(\varpi)=(\varpi,D^+(\gamma))_v$
is $+1$ or $-1$ according as the unit $D^+(\gamma)$ is a square or not;
both occur.

(d) \emph{$F=\Q_3$, $E=\Q_3(\sqrt3)$.} Let $h=\begin{pmatrix}0&3\\1&0
\end{pmatrix}$, which is $\sqrt3\in E^\times$, of determinant $-3$. Then
$\mathrm{Ad}(h)(-1,\zeta)=(-1,(-3,-1)_3\zeta)=(-1,-\zeta)$. By continuity
the elements of $\overline{T_3}$ close to a lift of $-1$ are conjugate
under $G_1(\Q_3)$, indeed under the centralizer $E^\times$, to their
negatives; here $D^+(\gamma)\in-3\,\Q_3^{\times2}$ and
$(-3,-3)_3=-1$.

(e) \emph{A rational example at $2$ and $\infty$.} Let
\[
\gamma=\begin{pmatrix}3/5&-4/5\\4/5&3/5\end{pmatrix},\qquad
h=\begin{pmatrix}1&-1\\1&1\end{pmatrix}.
\]
Both lie in $\Q[w]$, so they commute; $\det h=2$ and $D^+(\gamma)=16/5$.
Over $\Q_2$, $(2,16/5)_2=(2,5)_2=-1$: ordinary conjugation by $h$ exchanges
the two lifts of $\gamma$, and $h\star x=x$ for both. Over $\R$ the sign is
$(2,16/5)_\infty=+1$, and over $\Q_5$ it is $(2,5)_5=-1$, in agreement with
the product formula. (This example is due to the referee.)
\end{example}

\subsection{Stable conjugacy}\label{sec:stable-conj}

Two elements of $SL(2,F)$ are \emph{stably conjugate} if they are
conjugate in $SL(2,\bar F)$; for semisimple elements this is the same as
conjugacy under $G_1(F)$~\cite{LL}. On the cover, conjugacy under
$G_1(F)$ is too coarse, by Example~\ref{ex:eta}: it can identify the two
lifts of an element, on which every genuine function takes opposite
values. The remedy is to twist the action by the sign of
Lemma~\ref{lem:eta}.

Let $\Gb^\flat$ be the set of $x\in\Gb$ whose image $p(x)$ has trace
different from $-2$; it is open, dense, invariant under every
$\mathrm{Ad}(h)$, and contains the regular semisimple elements. For
$x\in\Gb^\flat$ write $\eta_x=\eta_{p(x)}$, and for $h\in G_1(F)$ put
\begin{equation}\label{eq:star}
h\star x=\eta_x(\det h)\cdot\mathrm{Ad}(h)x,\qquad
\eta_x(\det h)=\bigl(\det h,\ \det(1+p(x))\bigr)_v\in\mu_2 .
\end{equation}
Since $\eta_x$ depends only on the trace of $p(x)$, which
$\mathrm{Ad}(h)$ preserves, $(h_1h_2)\star x=h_1\star(h_2\star x)$: this
is an action of $G_1(F)$ on $\Gb^\flat$. Scalars act trivially, and
$h\in SL(2,F)$ acts by conjugation. Two elements of $\Gb^\flat$ are
called \emph{stably conjugate} if they lie in the same orbit of this
action. For a function $f$ on $\Gb$ put
\[
f^{[h]}(x)=f(h\star x)=\eta_x(\det h)\,f^h(x)\qquad(x\in\Gb^\flat),
\]
the last equality for anti-genuine $f$. The complement of $\Gb^\flat$ has
measure $0$, so $f^{[h]}$ is defined almost everywhere; it is smooth on
$\Gb^\flat$ but in general not across $\operatorname{tr}=-2$. A locally
integrable class function $\Theta$ on $\Gb$, and the invariant
distribution $f\mapsto\int f\,\Theta$ it defines, are called
\emph{stable} if $\Theta(h\star x)=\Theta(x)$ for all $h\in G_1(F)$ and
almost all $x$. Since $x\mapsto h\star x$ preserves the Haar measure,
this means $\int f^{[h]}\Theta=\int f\,\Theta$.

By Lemma~\ref{lem:eta}, a regular semisimple $x$ over $\gamma$ is fixed
by $h\star$ for every $h$ in the centralizer $E^\times$ of $\gamma$. So
the stable class of $x$ maps bijectively onto the $G_1(F)$-conjugacy
class of $\gamma$; it never contains the other lift of $\gamma$; and it
is the union of the conjugacy classes of the $\delta_a\star x$, $a\in
F^\times/N E^\times$: two classes if $E$ is a field, and one if $\gamma$
is hyperbolic.

That this is the right notion is shown by the following function, which
separates the two lifts and is constant on stable classes. Let
$\gamma_v(a)$ be the Weil index of~\eqref{eq:weil-index}, for the fixed
character $\psi=\psi_v$.

\begin{lemma}\label{lem:Delta}
Let $\Gb_{\mathrm{reg}}$ be the set of regular semisimple elements of
$\Gb$.
\begin{enumerate}[label=(\alph*)]
\item There is exactly one function $\Delta_\psi\colon
  \Gb_{\mathrm{reg}}\to\mu_4$ that is invariant under conjugation by
  $\Gb$ and satisfies
  \begin{equation}\label{eq:Delta}
  \Delta_\psi(g,\zeta)=\zeta\,\bigl(c,D^+(g)\bigr)_v\,
  \frac{\gamma_v(D^+(g))}{\gamma_v(1)}\qquad\text{for }
  g=\begin{pmatrix}\alpha&b\\c&d\end{pmatrix},\ c\ne0.
  \end{equation}
  It is genuine: $\Delta_\psi(x\zeta)=\zeta\,\Delta_\psi(x)$.
\item $\Delta_\psi(\mathrm{Ad}(h)x)=\eta_x(\det h)\,\Delta_\psi(x)$ for
  $h\in G_1(F)$; that is, $\Delta_\psi(h\star x)=\Delta_\psi(x)$. So two
  elements of $\Gb_{\mathrm{reg}}$ are stably conjugate if and only if
  their images in $SL(2,F)$ are conjugate under $G_1(F)$ and
  $\Delta_\psi$ takes the same value on them.
\item On the split torus, $\Delta_\psi(a(t),\zeta)=\zeta\,\gamma_v(1)/
  \gamma_v(t)=\chi_{\psi_{-1},v}(a(t),\zeta)$, the complex conjugate of
  Weil's character $\chi_{\psi,v}$ (Lemma~\ref{lem:weil}).
\item For $a\in F^\times$, $\Delta_{\psi_a}=\eta_x(a)\,\Delta_\psi$, where
  $\psi_a(u)=\psi(au)$.
\end{enumerate}
\end{lemma}

\begin{proof}
Write $\Delta$ for the right side of~\eqref{eq:Delta}, a function on the
set $X$ of $(g,\zeta)$ with $c\ne0$ and $\operatorname{tr}g\ne-2$; and
$D=D^+(g)$, which is invariant under conjugation. We show that
$\Delta(yxy^{-1})=\Delta(x)$ when $x$ and $yxy^{-1}$ lie in $X$.

\emph{Conjugation by $(n(u),1)$.} The lower left entry is unchanged, and
by~\eqref{eq:inverse} the sign is $\beta(n(u),g)\beta(n(u)g,n(-u))
\beta(n(u),n(-u))=(c,1)_v(1,c)_v=1$.

\emph{Conjugation by $(a(s),1)$.} The lower left entry becomes
$c/s^2$, and the sign is
\[
\beta(a(s),g)\,\beta(a(s)g,a(s)^{-1})\,\beta(a(s),a(s)^{-1})
=(c,s)_v\cdot(s,cs)_v\cdot(s,-1)_v=1 .
\]
So $\Delta$ is unchanged.

\emph{Conjugation by $(w,1)$, when $b\ne0$.} Here $wgw^{-1}=
\begin{pmatrix}d&-c\\-b&\alpha\end{pmatrix}$. Suppose $\alpha\ne0$.
Then $x(wg)=\alpha$, and the sign is $\beta(w,g)\beta(wg,w^{-1})
\beta(w,w^{-1})=(\alpha,\alpha/c)_v(-b/\alpha,b)_v=
(\alpha,-c)_v(\alpha,b)_v=(\alpha,-b c)_v$. So invariance amounts
to $(-b,D)_v(\alpha,-b c)_v=(c,D)_v$, that is, to
$(-b c,\ \alpha D)_v=1$. Put $X_0=-b c=1-\alpha d$ and
$Y_0=\alpha+1$. Then $\alpha D=\alpha^2+\alpha d+2\alpha=Y_0^2-X_0$, and
$(X_0,Y_0^2-X_0)_v=1$ because $Y_0^2-X_0$ is the norm of $Y_0+\sqrt{X_0}$.
If $\alpha=0$, then $b c=-1$, $x(wg)=b$, the sign is
$(b,b/c)_v(-1,b)_v=(b,-b^2)_v(-1,b)_v=1$, and the
identity to be checked is $(-b,D)_v=(c,D)_v$, i.e.\ $(-b c,D)_v=1$,
which holds because $-b c=1$.

Now let $y\in\Gb$ be arbitrary with $x,yxy^{-1}\in X$. If $p(y)\in B$,
then $y$ is a product of elements of the first two kinds and of $\mu_2$.
Otherwise $p(y)=b_1wn$ with $b_1\in B(F)$ and $n\in N(F)$ (Bruhat). The
element $x'=(n,1)x(n,1)^{-1}$ is in $X$, and its upper right entry is not
$0$: otherwise $(w,1)x'(w,1)^{-1}$, hence $yxy^{-1}$, would be upper
triangular. So the three steps apply in turn. Finally, every regular
semisimple element is conjugate to an element of $X$ (if $c=0\ne b$,
conjugate by $w$; if $c=b=0$, first by $n(1)$), which gives existence
and uniqueness of $\Delta_\psi$. The values lie in $\mu_4$, because $(\gamma_v(D)/\gamma_v(1))^2=(D,-1)_v$ by Lemma~\ref{lem:weil}(a), and
genuineness is clear.

(b) For $h\in SL(2,F)$ this is~(a). For $h=\delta_a$ and $x=(g,\zeta)\in
X$, \eqref{eq:Ad-delta} gives $\mathrm{Ad}(\delta_a)x=(\delta_ag
\delta_a^{-1},\zeta)$ with lower left entry $c/a$, so
$\Delta(\mathrm{Ad}(\delta_a)x)=(a,D)_v\Delta(x)$. Since $G_1(F)$ is
generated by $Z(F)$, $SL(2,F)$ and the $\delta_a$, the formula follows.
For the last statement, if $p(x')=hp(x)h^{-1}$, then $x'=\pm\,h\star x$,
and $\Delta_\psi(x')=\pm\Delta_\psi(x)$ accordingly.

(c) $(a(t),1)$ is conjugate, by $(n(1),1)$ and then $(w,1)$, to
$(g_2,(t^2-1,t)_v)$ with $g_2=\begin{pmatrix}t^{-1}&0\\t-t^{-1}&t
\end{pmatrix}$: the first conjugation has sign $1$, and the second has
sign $\beta(w,g_1)\beta(wg_1,w^{-1})=(1-t^{-2},t)_v=(t^2-1,t)_v$, where
$g_1=n(1)a(t)n(-1)$. Here
$D=(t+1)^2/t$, so $(c,D)_v=((t^2-1)/t,t)_v=(t^2-1,t)_v(t,-1)_v$ and
$\Delta_\psi(a(t),1)=(t,-1)_v\gamma_v(t)/\gamma_v(1)$. By
Lemma~\ref{lem:weil}(a), $(\gamma_v(t)/\gamma_v(1))^2=(t,-1)_v$, so this
is $\gamma_v(1)/\gamma_v(t)$; and $\chi_{\psi_{-1}}=\chi_\psi\cdot
(-1,\cdot)_v$ by Lemma~\ref{lem:weil}(d).

(d) $\gamma_{\psi_a,v}(D)/\gamma_{\psi_a,v}(1)=\gamma_v(aD)/\gamma_v(a)=
(a,D)_v\,\gamma_v(D)/\gamma_v(1)$, by Lemma~\ref{lem:weil}(a).
\end{proof}

So stable conjugacy can be described without $GL(2)$: \emph{$x$ and $x'$
are stably conjugate when $p(x)$, $p(x')$ are stably conjugate in
$SL(2,F)$ and $\Delta_\psi(x)=\Delta_\psi(x')$}; by~(d) the condition does
not depend on $\psi$. This is the definition of Adams~\cite{Ad} and
Li~\cite[\S5]{Li11}, who use, in place of $\Delta_\psi$, the character of
the difference of the two halves of the Weil representation. The two
functions agree up to a positive factor:

\begin{lemma}\label{lem:weil-kernel}
Let $\omega=\omega_{\psi'}$ be the Weil representation
of~\eqref{eq:weil-rep}, and $\gamma'$ the Weil index of $\psi'$. Then
$\omega(w,1)\phi(x)=\gamma'(1)^{-1}\int\phi(y)\psi'(-2xy)\,dy$, and for
$g=\begin{pmatrix}\alpha&b\\c&d\end{pmatrix}$ with $c\ne0$,
\begin{equation}\label{eq:weil-kernel}
\omega(g,1)\phi(x)=\gamma'(c)^{-1}\,|c|_v^{-1/2}\int_{F}\psi'\Bigl(
\frac{\alpha x^2-2xy+dy^2}{c}\Bigr)\phi(y)\,dy,
\end{equation}
with the measure that is self-dual for $\psi'(2xy)$. Moreover
$\omega_{\psi'}\circ\mathrm{Ad}(\delta_a)=\omega_{\psi'_a}$ for $a\in
F^\times$, as representations on $\mathcal S(F)$; in particular
$r^\pm_{\psi'}\circ\mathrm{Ad}(\delta_a)=r^\pm_{\psi'_a}$.
\end{lemma}

\begin{proof}
In the proof of Lemma~\ref{lem:weil-rep} it is shown that
$\omega(w(t),1)\phi(x)=\gamma'(t)|t|^{1/2}(\mathcal F\phi)(tx)$ for
$w(t)=\begin{pmatrix}0&t\\-t^{-1}&0\end{pmatrix}$; for $t=-1$ this is the
first formula, since $w(-1)=w$ and $\gamma'(-1)=\gamma'(1)^{-1}$. For
$c\ne0$, $g=n(\alpha/c)\,a(c^{-1})\,w\,n(d/c)$, and the same holds for
the lifts $(\cdot,1)$, because the three values of Kubota's cocycle are
$(c,1)_v$, $(1,c)_v$ and $(1,c)_v$. Apply~\eqref{eq:weil-rep},
Lemma~\ref{lem:weil-rep} with $t=c^{-1}$, and $\gamma'(-1)\gamma'(1)=1$.
For the last statement it suffices to compare the two representations on
$(n(u),1)$ and $(w,1)$, which generate $\Gb$ together with $\mu_2$. By
Lemma~\ref{lem:Ad}(d), $\omega_{\psi'}(\mathrm{Ad}(\delta_a)(n(u),1))$ is
multiplication by $\psi'(aux^2)$, and
$\omega_{\psi'}(\mathrm{Ad}(\delta_a)(w,1))\phi(x)=\frac{\gamma'(1)}
{\gamma'(a)}|a|^{1/2}\cdot\gamma'(1)^{-1}\int\phi(y)\psi'(-2axy)\,dy$. This
is $\omega_{\psi'_a}(w,1)\phi(x)$: the Weil index of $\psi'_a$ at $1$ is
$\gamma'(a)$, and the measure self-dual for $\psi'_a(2xy)$ is $|a|^{1/2}$
times the one for $\psi'(2xy)$.
\end{proof}

As a check on the constant in~\eqref{eq:weil-kernel}: composing the
kernels of $(g_1,1)$ and $(g_2,1)$ by a Gaussian integral gives the
kernel of $g_1g_2$ times $\gamma'(c_3/c_1c_2)\gamma'(c_3)/\gamma'(c_1)
\gamma'(c_2)$, where $c_3$ is the lower left entry of $g_1g_2$; by
Lemma~\ref{lem:weil}(a) this is $(c_3/c_1,c_3/c_2)_v=\beta(g_1,g_2)$. So
\eqref{eq:weil-kernel} is a genuine representation of the cover defined
by Kubota's cocycle, with no further constant.

\begin{proposition}\label{prop:weil-char}
Let $\Theta^\pm_{\psi}$ be the characters of $r^\pm_{\psi}$. On
$\Gb_{\mathrm{reg}}$ the distribution $\Theta^+_\psi-\Theta^-_\psi$ is the
function
\begin{equation}\label{eq:weil-char}
\Theta^+_\psi(x)-\Theta^-_\psi(x)=\frac{\Delta_\psi(x)}{|D^+(x)|_v^{1/2}} ;
\end{equation}
that is, $\operatorname{tr}r^+_\psi(f)-\operatorname{tr}r^-_\psi(f)=\int
f(x)\Delta_\psi(x)|D^+(x)|_v^{-1/2}dx$ for $f\in C_c^\infty(\Gb)$ with
support in $\Gb_{\mathrm{reg}}$. Hence $\Theta^+_\psi-\Theta^-_\psi$ is
invariant under $\star$ on $\Gb_{\mathrm{reg}}$, it does not vanish there,
and stable conjugacy as defined above coincides with stable conjugacy in
the sense of~\cite{Ad,Li11}. The right side of~\eqref{eq:weil-char} is
locally integrable on $\Gb$, and by the regularity theorem for the
character of the Weil
representation~\cite{Mak,Th08} the formula holds for all $f\in
C_c^\infty(\Gb)$: $\Theta^+_\psi-\Theta^-_\psi$ is a stable character.
\end{proposition}

\begin{proof}
Let $P\phi(x)=\phi(-x)$, so that $\operatorname{tr}r^+_\psi(f)-
\operatorname{tr}r^-_\psi(f)=\operatorname{tr}(\omega_\psi(f)P)$. Both
sides of the identity to be proved are unchanged when $f$ is replaced by a
conjugate $f^y$. Let $X$ be the set of the proof of
Lemma~\ref{lem:Delta}; it is open, and every element of
$\Gb_{\mathrm{reg}}$ is conjugate to an element of it. So each point of the
support of $f$ has an open neighborhood of the form $yUy^{-1}$ with
$U\subset X$ open. Finitely many of these cover the support, and a
partition of unity subordinate to them writes $f$ as a finite sum of
functions $f_i\in C_c^\infty(\Gb)$ with $f_i^{y_i}$ supported in $X$. We
may therefore assume that $f$ has compact support in $X$. Then $c$ and
$D=D^+(g)$ are bounded away from $0$ and $\infty$ on the support of $f$.

For $x=(g,\zeta)\in X$ let $k_x(\xi,\eta)$ be $\zeta$ times the kernel
of~\eqref{eq:weil-kernel}, so that $|k_x|=|c|_v^{-1/2}$ and
\[
k_x(\xi,-\xi)=\zeta\,\gamma_v(c)^{-1}|c|_v^{-1/2}\,\psi\bigl(D\xi^2/c\bigr).
\]
By Fubini's theorem $\omega_\psi(f)$ has the kernel $K_f(\xi,\eta)=\int
f(x)k_x(\xi,\eta)\,dx$, which is continuous and bounded.

\emph{Let $v$ be finite.} Then $f$ is bi-invariant under the image $K'$ of
some $K_m$, and $\omega_\psi(f)=\omega_\psi(e')\omega_\psi(f)\omega_\psi(e')$
for the idempotent $e'$ of $K'$. The space $V'$ of the vectors of
$L^2(F)$ fixed by $K'$ is finite-dimensional and consists of
Schwartz--Bruhat functions, by a two-lattice argument. The elements $(n(u),1)$, $u\in\varpi^m\OO$, act by multiplication
by $\psi(u\xi^2)$, so a fixed vector is supported in the bounded set
$L_1=\{\xi\mid\psi(u\xi^2)=1\text{ for all }u\in\varpi^m\OO\}$. The
conjugates by $(w,1)$ of the $(n(u),1)$ with $u\in2\varpi^m\OO$ lie in $K'$
(they lie over $K_m$, and the sign is a character of $u\in\varpi^m\OO$,
hence trivial on $2\varpi^m\OO$), and they act in the same way on the
Fourier transform, by~\eqref{eq:weil-rep}; so the Fourier
transform of a fixed vector is supported in a bounded set as well, and the
vector is constant on the cosets of a lattice $L_2$. Functions supported in a
bounded set and constant modulo a lattice form a finite-dimensional space.
So $\omega_\psi(f)$ has finite rank: $\omega_\psi(f)\phi=\sum_i\langle
\phi,w_i\rangle u_i$ with finitely many $u_i,w_i\in V'$ (the functionals
are of this form because $\omega_\psi(e')$ is the orthogonal projection
onto $V'$). Its kernel is $\sum_iu_i(\xi)\overline{w_i(\eta)}$, a
Schwartz--Bruhat function, which equals the continuous function $K_f$,
and $\operatorname{tr}(\omega_\psi(f)P)=\sum_i\langle Pu_i,w_i\rangle=
\int_FK_f(\xi,-\xi)\,d\xi$. Hence
\[
\operatorname{tr}(\omega_\psi(f)P)=\lim_{M\to\infty}\int f(x)\Bigl(
\int_{\varpi^{-M}\OO}k_x(\xi,-\xi)\,d\xi\Bigr)dx .
\]
By Lemma~\shellref below, applied to $a=D/c$, the inner integral
does not depend on $M$ for $M\ge M_0$, with $M_0$ uniform on the support
of $f$, and by~\eqref{eq:weil-index} (the measure being $|2|_v^{1/2}$ times
the one there) its value is $\zeta$ times
\[
\gamma_v(c)^{-1}|c|_v^{-1/2}\int_F\psi\Bigl(\frac{D\,\xi^2}{c}\Bigr)d\xi
=\frac{\gamma_v(cD)}{\gamma_v(c)}\,|D|_v^{-1/2}
=(c,D)_v\,\frac{\gamma_v(D)}{\gamma_v(1)}\,|D|_v^{-1/2};
\]
the last step is Lemma~\ref{lem:weil}(a). This is~\eqref{eq:weil-char}.

\emph{Let $v$ be archimedean.} We show that $K_f$ is a Schwartz function
on $F\times F$. On $X$ use the coordinates $(\alpha,c,d)$. The section
$g\mapsto(g,1)$ is continuous on each connected component of $\{c\ne0\}$:
$g=n(\alpha/c)\,a(-c^{-1})\,w\,n(d/c)$, the sections of
Lemma~\ref{lem:local}(i) are continuous on the factors, and the cocycle
of the factors is a Hilbert symbol in $c$, which is locally constant. The Haar
measure is a smooth
multiple of $d\alpha\,dc\,dd$, and $\gamma_v(c)$ is locally constant. So
\[
K_f(\xi,\eta)=\int_{F^\times}\widehat\Phi_c\Bigl(\frac{\xi^2}c,\frac{\eta^2}c
\Bigr)\,\psi\Bigl(-\frac{2\xi\eta}c\Bigr)\,dc,
\]
where $\Phi(\alpha,c,d)$ is smooth with compact support in $\{c\ne0\}$ and
$\widehat\Phi_c$ is its Fourier transform in $(\alpha,d)$. The functions
$\widehat\Phi_c$ and their derivatives in $c$ are Schwartz functions,
uniformly for $c$ in the compact set in question, and $|\xi|^2+|\eta|^2$
is bounded by a constant times $|\xi^2/c|+|\eta^2/c|$ there. So $K_f$ decays
faster than any power of $1+|\xi|+|\eta|$, and so do its derivatives, which
are integrals of the same kind with polynomial factors. An integral
operator on $L^2(F)$ with a Schwartz kernel is of trace class, and its
trace is the integral of the kernel over the diagonal.
So $\operatorname{tr}(\omega_\psi(f)P)=\int K_f(\xi,-\xi)\,d\xi=\lim_{\epsilon
\to0}\int f(x)\bigl(\int
k_x(\xi,-\xi)e^{-\epsilon|\xi|^2}d\xi\bigr)dx$. The inner integrals are
Gaussian; they are bounded by a constant times $|D|_v^{-1/2}$, uniformly
in $\epsilon$, and converge to the value above. Dominated convergence
gives~\eqref{eq:weil-char}.

The remaining statements follow from Lemma~\ref{lem:Delta}(b) and
$\Delta_\psi\ne0$. Local integrability: $|D^+|_v^{-1/2}$ is the inverse
square root of the function $2+\operatorname{tr}$, which vanishes to first
order on the smooth part of its zero set and is a nondegenerate quadratic
form in three variables at $-1$.
\end{proof}

\begin{lemmaprime}
Let $F$ be non-archimedean, let $\psi$ be trivial on $\varpi^{n_0}\OO$ and
not on $\varpi^{n_0-1}\OO$, and let $e=\mathrm{ord}(2)$. For $a\in F^\times$
and $j\in\Z$ with
\[
2j\ \ge\ \mathrm{ord}(a)-n_0+2+2e,
\]
the integral of $\psi(a\xi^2)$ over the shell $\{\mathrm{ord}(\xi)=-j\}$
vanishes. Consequently $\int_{\varpi^{-M}\OO}\psi(a\xi^2)\,d\xi$ does not
depend on $M$ once $2(M+1)\ge\mathrm{ord}(a)-n_0+2+2e$, and the bound is
uniform for $a$ in a compact subset of $F^\times$.
\end{lemmaprime}

\begin{proof}
Put $m=\mathrm{ord}(a)-2j$, the order of $a\xi^2$ on the shell, and
$r=n_0-1-m-e$. The hypothesis is $m\le n_0-2-2e$; it gives $r\ge1+e\ge1$
and $m+2r\ge n_0$. For $u\in\varpi^r\OO$ the substitution
$\xi\mapsto\xi(1+u)$ preserves the shell and the measure, and
\[
\psi\bigl(a\xi^2(1+u)^2\bigr)=\psi(a\xi^2)\,\psi(2a\xi^2u)\,\psi(a\xi^2u^2)
=\psi(a\xi^2)\,\psi(2a\xi^2u),
\]
because $\mathrm{ord}(a\xi^2u^2)\ge m+2r\ge n_0$. Averaging over
$u\in\varpi^r\OO$ multiplies the integrand by the average of the character
$u\mapsto\psi(2a\xi^2u)$ of $\varpi^r\OO$, which is not trivial because
$2a\xi^2\varpi^r\OO=\varpi^{m+e+r}\OO=\varpi^{n_0-1}\OO$. So the integral
vanishes.
\end{proof}

For $F=\Q_2$ and $\psi(x)=e^{2\pi i\{x\}}$ the lemma says that the
integrals over $2^{-M}\Z_2$ are constant from $M=1$ on for $a=1$ and from
$M=2$ on for $a=2$; the values are $1+i$ and $\sqrt2(1+i)$.

The same computation without $P$ gives the character of $\omega_\psi$
itself on the set where $D^-=\operatorname{tr}g-2\ne0$:
\begin{equation}\label{eq:weil-char-full}
\Theta_{\omega_\psi}(g,\zeta)=\zeta\,(c,D^-)_v\,\frac{\gamma_v(D^-)}
{\gamma_v(1)}\,|D^-|_v^{-1/2},\qquad c\ne0 .
\end{equation}
Two comparisons. First, $\det(g-1)=-D^-$ and $\gamma_v(1)\gamma_v(-D^-)=
(-1,D^-)_v\gamma_v(D^-)/\gamma_v(1)$, so~\eqref{eq:weil-char-full} is
$(-c,D^-)_v\cdot\gamma_v(1)\gamma_v(\det(g-1))|\det(g-1)|_v^{-1/2}$. This is
the formula of Thomas~\cite[Theorem~1C]{Th08}. That theorem is stated up
to a sign, so the comparison confirms the absolute value and the Weil
index and says nothing about the sign $(-c,D^-)_v$, which rests on the
computation above. Second,
$\omega_\psi(-1,1)=\gamma_v(1)^2P$ (Lemma~\ref{lem:weil-rep}), so
$\Theta^+_\psi(x)-\Theta^-_\psi(x)=\gamma_v(1)^{-2}\,\Theta_{\omega_\psi}
\bigl((-1,1)x\bigr)$; and indeed~\eqref{eq:weil-char-full} at
$(-1,1)(g,1)=(-g,(c,-1)_v)$ gives back~\eqref{eq:weil-char}, by
Lemma~\ref{lem:weil}(a) and $\gamma_v(1)^4=(-1,-1)_v$. The character of the
difference in~\cite[Theorem~D]{Th13} stands in the same relation to the
character of the Weil representation at $-g$; that theorem is written in
other coordinates on the metaplectic group and with a Weil index of a
quadratic form, and we have not verified an equality between its factors
and ours, nor matched the normalizations of the Weil index
in~\cite{Th08,Th13} with ours. Lemma~\ref{lem:Delta}, on which
the rest of the section depends, is algebraic and uses none of this
(Remark~\ref{rem:sec7-status}).

\begin{example}\label{ex:Delta-real}
Let $F=\R$. For $0<\theta<\pi$ the element $\tilde k_\theta=(k_\theta,1)$
has $c=-\sin\theta<0$ and $D^+=4\cos^2\frac\theta2>0$, so
$\Delta_\psi(\tilde k_\theta)=1$ and~\eqref{eq:weil-char} gives
$\Theta^+-\Theta^-=1/(2\cos\frac\theta2)$. Directly: if $r^+_\psi$ has
lowest weight $\frac12$, its weights are $\frac12+2j$ and those of
$r^-_\psi$ are $\frac32+2j$, $j\ge0$, and
\[
\sum_{j\ge0}e^{i(\frac12+2j)\theta}-\sum_{j\ge0}e^{i(\frac32+2j)\theta}
=\frac{e^{i\theta/2}}{1+e^{i\theta}}=\frac1{2\cos\frac\theta2}
\]
(Abel summation). The function is even in $\theta$, in agreement with
Example~\ref{ex:eta}(b): $\tilde k_\theta$ and $\tilde k_{-\theta}=
\mathrm{Ad}(\delta_{-1})\tilde k_\theta$ are stably conjugate, and they
are not conjugate. On the hyperbolic elements over $a(t)$, $t<0$, the
stable class is the conjugacy class, and $\mathrm{Ad}(\delta_{-1})$
exchanges the two lifts.
\end{example}

\begin{example}\label{ex:Delta-finite}
Let $v$ be finite and non-dyadic, $\psi$ unramified, and $x=(g,\kappa(g))
\in K^*$ with $c$ and $D^+(g)$ units. Then $\Delta_\psi(x)=1$ and
$\Delta_{\psi_\varpi}(x)=(\varpi,D^+(g))_v$, the quadratic residue symbol
of $\det(1+g)$. The second value can be seen in the Weil representation
of $SL(2,\mathbb F_q)$ on the functions on $\mathbb F_q$, which is the
space of vectors of $\omega_{\psi_\varpi}$ fixed by $K_1$ (the functions on
$\varpi^{-1}\OO/\OO$): there the trace of $\omega(g)P$ is
$\bigl(\frac{\det(1+g)}{q}\bigr)$ for every $g$ with $\det(1+g)\ne0$. The
script verifies this for $q=3,5,7,11$ by computing the matrices. (The
character of $\omega_{\psi_\varpi}$ also involves the vectors of higher
level, which we have not computed separately.)
\end{example}

\paragraph{The dual group.} The unramified genuine principal series of
$\Gb$ are parametrized by $\chi_v(\varpi)\in\C^\times$ modulo inversion,
through $\pi_v=\chi_{\psi,v}\cdot(\chi_v\circ p)$, and the intertwining
operator has the normalizing factor $L_v(2\lambda,\chi_v^2)/L_v(2\lambda+1,
\chi_v^2)$ (Lemma~\ref{lem:intertwining}). This is the normalizing factor
of the principal series of $PGL(2,F)$ induced from $\diag(t_1,t_2)\mapsto
\chi_v(t_1/t_2)$, whose Satake parameter is $\diag(\chi_v(\varpi),
\chi_v(\varpi)^{-1})\in SL(2,\C)$. So the dual group of the double cover
of $SL(2)$ is $SL(2,\C)$, the dual group of $SO(3)=PGL(2)$, and not
$PGL(2,\C)$, the dual group of $SL(2)$~\cite{Sav,Li11,GGW,We18}. The
elliptic endoscopic data correspond to the two central elements $\pm1$
of $SL(2,\C)$. This is why the endoscopic groups are $SO(3)$ twice, why
the translation by $-1$ appears below, and why no tori occur.

\subsection{Stable and \texorpdfstring{$\kappa$}{kappa}-orbital
integrals}\label{sec:kappa}

We now sort the orbital integrals in a $G_1(F)$-conjugacy class by the
characters of $F^\times/N E^\times$. For $SL(2)$ this gives the stable
orbital integral and one $\kappa$-orbital integral~\cite{LL}. On the
cover there are again two combinations, and the second is a stable
orbital integral at another point.

Fix Haar measures on the centralizers so that they are compatible under
every $\mathrm{Ad}(h)$, $h\in G_1(F)$ (Section~\ref{sec:measures-local}).
For $x\in\Gb_{\mathrm{reg}}$ over $\gamma$, with centralizer $T_x=T_\gamma$
in $SL(2,F)$, and $f\in C_c^\infty(\Gb)_-$, let
\[
J_G(x,f)=|D(\gamma)|_v^{1/2}\int_{T_x\backslash SL(2,F)}f(y^{-1}xy)\,dy,
\qquad D(\gamma)=(\operatorname{tr}\gamma)^2-4,
\]
as in~\eqref{eq:JG}. Then $J_G(x,f^h)=J_G(\mathrm{Ad}(h)x,f)$ and
$J_G(x,f^{[h]})=J_G(h\star x,f)$ for $h\in G_1(F)$. Let $E=F[\gamma]$,
and define
\begin{equation}\label{eq:Jst}
J^{\mathrm{st}}_G(x,f)=\sum_{a\in F^\times/N E^\times}J_G(\delta_a\star x,f),
\qquad
J^{\kappa}_G(x,f)=\sum_{a\in F^\times/N E^\times}\kappa_E(a)\,
J_G(\delta_a\star x,f).
\end{equation}
The terms depend only on $a$ modulo $N E^\times$, because $\delta_{aN(e)}
\star x$ is conjugate to $\delta_a\star(e\star x)=\delta_a\star x$ for
$e\in E^\times$. If $\gamma$ is hyperbolic, both are $J_G(x,f)$.

\begin{lemma}\label{lem:kappa-orbital}
Let $x\in\Gb_{\mathrm{reg}}$ lie over $\gamma$, $E=F[\gamma]$, and
$f\in C_c^\infty(\Gb)_-$.
\begin{enumerate}[label=(\alph*)]
\item $J^{\mathrm{st}}_G(h\star x,f)=J^{\mathrm{st}}_G(x,f)$ and
  $J^\kappa_G(h\star x,f)=\kappa_E(\det h)J^\kappa_G(x,f)$ for $h\in
  G_1(F)$. Equivalently, $J^{\mathrm{st}}_G(x,f^{[h]})=J^{\mathrm{st}}_G(x,
  f)$: the distribution $J^{\mathrm{st}}_G(x,\cdot)$ is stable.
\item For a character $\kappa'$ of $F^\times/F^{\times2}$,
  \[
  \sum_{a\in F^\times/F^{\times2}}\kappa'(a)\,J_G(x,f^{\delta_a})
  =[N E^\times:F^{\times2}]\cdot\begin{cases}
  J^{\mathrm{st}}_G(x,f),&\kappa'=\eta_\gamma,\\
  J^{\kappa}_G(x,f),&\kappa'=\eta_\gamma\kappa_E,\\
  0,&\text{otherwise.}\end{cases}
  \]
  In particular the average of $J_G(x,f^h)$ over $h\in G_1(F)/Z(F)SL(2,F)$
  vanishes unless $D^+(\gamma)\in F^{\times2}\cup\theta F^{\times2}$; for
  $\gamma=a(t)$, unless $t$ is a square.
\item Let $E$ be a field, and let $\tilde z$ be a lift of $-1$. Then
  \[
  J^\kappa_G(x,f_{\tilde z})=J^{\mathrm{st}}_G(\tilde zx,f),
  \qquad\text{that is,}\qquad
  J^\kappa_G(x,f)=J^{\mathrm{st}}_G(\tilde zx,f_{\tilde z^{-1}}).
  \]
\item Let $v$ be finite and non-dyadic, $f=e_v$, $x\in K^*$ and
  $|D(\gamma)|_v=1$, and let $T_\gamma\cap K$ have volume $1$. Then
  $J^{\mathrm{st}}_G(x,e_v)=J^\kappa_G(x,e_v)=J_G(x,e_v)=1$.
\end{enumerate}
\end{lemma}

\begin{proof}
(a) $\delta_a\star(h\star x)=(\delta_ah)\star x$. Write
$\delta_ah=\delta_{a\det h}h_0$ with $h_0\in SL(2,F)$; then $(\delta_ah)
\star x$ is conjugate to $\delta_{a\det h}\star x$, by
$\delta_{a\det h}h_0\delta_{a\det h}^{-1}$. So the terms of~\eqref{eq:Jst}
are permuted by $a\mapsto a\det h$. The reformulation is
$J_G(y,f^{[h]})=J_G(h\star y,f)$.

(b) $J_G(x,f^{\delta_a})=J_G(\mathrm{Ad}(\delta_a)x,f)=\eta_\gamma(a)
J_G(\delta_a\star x,f)$, since $f$ is anti-genuine. So the left side is
$\sum_a\kappa'(a)\eta_\gamma(a)J_G(\delta_a\star x,f)$, and the summand
$J_G(\delta_a\star x,f)$ depends only on $a$ modulo $N E^\times$. The sum
of $\kappa'\eta_\gamma$ over a coset of $N E^\times/F^{\times2}$ is $0$
unless $\kappa'\eta_\gamma$ is trivial on $N E^\times$, that is, unless
$\kappa'\eta_\gamma\in\{1,\kappa_E\}$. For $\kappa'=1$ the condition is
$\eta_\gamma\in\{1,\kappa_E\}$.

(c) Write $\gamma=\begin{pmatrix}x_0&y_0\theta\\y_0&x_0\end{pmatrix}$ in a
suitable basis, $x_0^2-\theta y_0^2=1$. Then
$D^+(\gamma)D^+(-\gamma)=4-(2x_0)^2=-4\theta y_0^2$, so
\begin{equation}\label{eq:eta-minus}
\eta_\gamma\,\eta_{-\gamma}=(\cdot,-\theta)_v=\kappa_E\cdot(\cdot,-1)_v .
\end{equation}
By Lemma~\ref{lem:eta}, $\mathrm{Ad}(\delta_a)\tilde z=(a,-1)_v\tilde z$.
Hence
\[
\delta_a\star(\tilde zx)=\eta_{-\gamma}(a)\,(a,-1)_v\,\tilde z\,
\mathrm{Ad}(\delta_a)x=\eta_{-\gamma}(a)(a,-1)_v\eta_\gamma(a)\cdot
\tilde z\,(\delta_a\star x)=\kappa_E(a)\,\tilde z\,(\delta_a\star x).
\]
Since $J_G(\tilde zy,f)=J_G(y,f_{\tilde z})$, $J_G(\varepsilon y,f)=
\varepsilon J_G(y,f)$ for $\varepsilon\in\mu_2$, and $D(-\gamma)=D(\gamma)$,
summing over $a$ gives $J^{\mathrm{st}}_G(\tilde zx,f)=J^\kappa_G(x,
f_{\tilde z})$.

(d) $J_G(x,e_v)=1$ is the unramified computation used for
Theorem~\ref{thm:fine}. If $E$ is a field, it is unramified, and
$\delta_\varpi\star x$ lies over $\delta_\varpi\gamma\delta_\varpi^{-1}$,
which has no $SL(2,F)$-conjugate in $K$ (proof of
Theorem~\ref{thm:coarse}(a)); so $J_G(\delta_\varpi\star x,e_v)=0$.
\end{proof}

Part~(b) explains a statement of the 1995 notes: for the double cover the
orbital integral of a genuine function at a lift of $a(t)$ does not
vanish when $t$ is not a square, but its average over the $G_1(F)$-orbit
does (Appendix~\ref{sec:appendix-1995}). Part~(c) is the form that the
relation between the endoscopic data $(1,0)$ and $(0,1)$ takes here: Li's
transfer factors for the two data differ by the translation by the
central element $-1$~\cite[\S5]{Li11}, and correspondingly the
$\kappa$-orbital integral of $f$ at $x$ is the stable orbital integral of
a translate of $f$ at $\tilde zx$. For $SL(2)$ there is no such
identity: there $J^\kappa$ is a new function on the torus $T_\gamma$, and
\cite{LL} apply Poisson summation on $T_\gamma$ to it.

By Lemma~\ref{lem:Delta}(b), $x\mapsto\Delta_\psi(x)^{-1}J^{\mathrm{st}}_G
(x,f)$ is invariant under $\mu_2$ and under $\star$, so it is a function
of the stable conjugacy class of $p(x)$ in $SL(2,F)$, that is, of the
trace. Li's transfer theorem~\cite{Li11} says that it is the stable
orbital integral of a function $f'$ on $SO(3,F)$ at the element with
eigenvalues $\epsilon,1,\epsilon^{-1}$, where $\epsilon^{\pm1}$ are the
eigenvalues of $p(x)$, up to a normalization of the transfer factor. We
do not use this.

\subsection{The elliptic regular terms}\label{sec:elliptic-stable}

Let $\gamma\in G(\Q)=SL(2,k)$ be regular elliptic, $E=k[\gamma]$ its
splitting field, $T_\gamma$ its centralizer and $\kappa_E=\prod_v\kappa_{E_v}$
the quadratic character of $k^\times\backslash\A_k^\times$ attached to
$E$. Everything in Sections~\ref{sec:Ad}--\ref{sec:kappa} has an adelic
version. Let
\[
\eta_\gamma(a)=\prod_v\bigl(a_v,D^+(\gamma)\bigr)_v,\qquad a\in\A_k^\times,
\]
a quadratic character of $\A_k^\times$ which is trivial on $k^\times$ by
the product formula~\ref{H4}, because $D^+(\gamma)=2+\operatorname{tr}
\gamma\in k^\times$. For $h\in G_1(\A_k)$ and $x$ in the orbit of
$\pm s(\gamma)$ under $\mathrm{Ad}(G_1(\A_k))$ put
$h\star x=\eta_\gamma(\det h)\,\mathrm{Ad}(h)x$; its component at $v$ is
$h_v\star x_v$ in the sense of~\eqref{eq:star}. So for $a\in k^\times$,
\begin{equation}\label{eq:rational-stable}
s(\delta_a\gamma\delta_a^{-1})=\mathrm{Ad}(\delta_a)s(\gamma)=\delta_a
\star s(\gamma)
\end{equation}
by Lemma~\ref{lem:Ad}(f): \emph{rational elements that are conjugate
under $G_1(k)$ are stably conjugate in $\GA$, at every place.} This is
what makes a stabilization possible, and it is the reason for the twist
in~\eqref{eq:star}: the untwisted local conjugates
$\mathrm{Ad}(\delta_a)\gamma_v$ are not coherent in $v$. (Compare
\cite[Lemme~4.1.7]{Li15}.)

For $f\in\mathcal H(\GA)$ and $x=a\star s(\gamma)$, $a\in\A_k^\times$
(we write $a\star$ for $\delta_a\star$), let $J_G(x,f)=\int f(y^{-1}xy)\,
dy$ over $T_x(\A)\backslash G(\A)$, with the Tamagawa measures on $G(\A)$
and on $T_x(\A)=\delta_aT_\gamma(\A)\delta_a^{-1}$; no factor $|D|^{1/2}$ is
needed, since $|D(\gamma)|_{\A_k}=1$. Put
\begin{equation}\label{eq:Jst-global}
J^{\mathrm{st}}_G(\gamma,f)=\sum_{a\in\A_k^\times/N\A_E^\times}J_G\bigl(a\star
s(\gamma),f\bigr),\qquad
J^{\kappa}_G(\gamma,f)=\sum_{a\in\A_k^\times/N\A_E^\times}\kappa_E(a)\,J_G\bigl(
a\star s(\gamma),f\bigr).
\end{equation}
Here $\A_k^\times/N\A_E^\times=\bigoplus_vk_v^\times/N E_v^\times$. By
Lemma~\ref{lem:kappa-orbital}(d) only finitely many $a$ contribute, and
for $f=\bigotimes f_v$ and local measures whose products are the Tamagawa
measures,
\[
J^{\mathrm{st}}_G(\gamma,f)=\prod_vJ^{\mathrm{st}}_G(\gamma_v,f_v),\qquad
J^{\kappa}_G(\gamma,f)=\prod_vJ^{\kappa}_G(\gamma_v,f_v),
\]
where $(\gamma_v)$ are local components of $s(\gamma)$; the products do
not depend on their choice, and almost all factors are $1$.

\begin{theorem}\label{thm:elliptic-stable}
For $f\in\mathcal H(\GA)$ let
\[
S_{\mathrm{ell}}(f)=\frac12\sum_{\gamma}\mathrm{vol}\bigl(T_\gamma(\Q)
\backslash T_\gamma(\A)\bigr)\,J^{\mathrm{st}}_G(\gamma,f),
\]
the sum over a set of representatives $\gamma$ of the $G_1(k)$-conjugacy
classes of regular elliptic elements of $SL(2,k)$. The sum is finite, and
\[
\sum_{\mathcal O\ \mathrm{elliptic}}J_{\mathcal O}(f)=S_{\mathrm{ell}}(f)
+S_{\mathrm{ell}}(f_-).
\]
More precisely, for each $\gamma$ the classes $\mathcal O$ contained in
the $G_1(k)$-conjugacy class of $\gamma$ contribute
\begin{equation}\label{eq:pre-stab}
\sum_{\mathcal O\subset G_1(k)\cdot\gamma}J_{\mathcal O}(f)=\frac12\,
\mathrm{vol}\bigl(T_\gamma(\Q)\backslash T_\gamma(\A)\bigr)\Bigl(
J^{\mathrm{st}}_G(\gamma,f)+J^{\kappa}_G(\gamma,f)\Bigr),
\qquad J^\kappa_G(\gamma,f)=J^{\mathrm{st}}_G(-\gamma,f_-).
\end{equation}
The distribution $S_{\mathrm{ell}}$ is stable: for $f=\bigotimes_vf_v$ it
is a finite sum of products of the stable local distributions
$J^{\mathrm{st}}_G(\gamma_v,\cdot)$ of Lemma~\ref{lem:kappa-orbital}(a).
\end{theorem}

\begin{proof}
The $SL(2,k)$-conjugacy classes in the $G_1(k)$-class of $\gamma$ are the
classes $\mathcal O_a$ of $\delta_a\gamma\delta_a^{-1}$, $a\in k^\times/
N E^\times$. By Proposition~\ref{prop:elliptic}
and~\eqref{eq:rational-stable}, $J_{\mathcal O_a}(f)=\mathrm{vol}\cdot
J_G(a\star s(\gamma),f)$, with the same volume for all $a$, the Tamagawa
measure being compatible with $\mathrm{Ad}(\delta_a)$. Let
$\Phi(a)=J_G(a\star s(\gamma),f)$ for $a\in\A_k^\times$. It depends only
on the image of $a$ in $\A_k^\times/N\A_E^\times$, and it vanishes outside
a finite set. By the Hasse norm theorem, $k^\times/N E^\times$ injects
into $\A_k^\times/N\A_E^\times$, and by class field theory its image is the
kernel of $\kappa_E$, of index $2$~\cite{CF}. Therefore
\[
\sum_{a\in k^\times/N E^\times}\Phi(a)=\frac12\sum_{a\in\A_k^\times/N\A_E^\times}
\bigl(1+\kappa_E(a)\bigr)\Phi(a)=\frac12\Bigl(J^{\mathrm{st}}_G(\gamma,f)+
J^\kappa_G(\gamma,f)\Bigr).
\]
For the second formula of~\eqref{eq:pre-stab} we may take
$f=\bigotimes f_v$. Choose local components $\gamma_v$ of $s(\gamma)$ and
$\tilde z_v$ of $s(-1)$. Then $(\tilde z_v\gamma_v)$ are local components
of $s(-1)s(\gamma)=s(-\gamma)$, and $\bigotimes_v(f_v)_{\tilde z_v^{-1}}=
f_{s(-1)^{-1}}=f_-$. By Lemma~\ref{lem:kappa-orbital}(c) at the places
where $E_v$ is a field (at the others both sides are $J_G$),
\[
J^\kappa_G(\gamma,f)=\prod_vJ^{\mathrm{st}}_G\bigl(\tilde z_v\gamma_v,
(f_v)_{\tilde z_v^{-1}}\bigr)=J^{\mathrm{st}}_G(-\gamma,f_-).
\]
Finally $\gamma\mapsto-\gamma$ is a bijection of the set of
$G_1(k)$-classes of regular elliptic elements, with $T_{-\gamma}=
T_\gamma$. Summing~\eqref{eq:pre-stab} over $\gamma$ gives the first
formula; the sums are finite by Theorem~\ref{thm:coarse}(a). Stability is
Lemma~\ref{lem:kappa-orbital}(a) at each place.
\end{proof}

For $SL(2)$ the first formula of~\eqref{eq:pre-stab} is the
pre-stabilization of~\cite[\S5]{LL}. The second has no analog there, and
it is the content of the theorem: the unstable half of the elliptic terms
of $f$ is the stable half of the elliptic terms of $f_-$. In the notation
of~\cite[Th\'eor\`eme~5.2.6]{Li15}, where the elliptic part of the trace
formula is $\sum_{(n',n'')}\iota(\tilde G,H)\,ST^H_{\mathrm{ell}}(f^H)$ with
$\iota=\frac12$ for both pairs when $n=1$, there are likewise two
summands with coefficient $\frac12$, those of $(1,0)$ and $(0,1)$. This is
a comparison of the decomposition and of the coefficients only. Li's terms
are stable distributions on $SO(3)$ applied to transfers of $f$, with
transfer factors; ours are distributions on the cover. Matching them
requires the transfer and its normalization, which are not carried out
here (Remark~\ref{rem:sec7-status}).

\begin{example}\label{ex:order4-stable}
Let $k=\Q$ and $\gamma=w$, of order $4$, with $E=\Q(i)$, $D^+(w)=2$ and
$\eta_w=\prod_v(\cdot,2)_v$. The $G_1(\Q)$-conjugacy class of $w$ is the
union of the $SL(2,\Q)$-classes of the $\delta_aw\delta_a^{-1}$,
$a\in\Q^\times/N E^\times$; there are infinitely many, and for a given $f$
finitely many contribute. The class $a=-1$ is that of
$w^{-1}=-w$. So $-w$ lies in the $G_1(\Q)$-class of $w$,
and~\eqref{eq:pre-stab} reads
\[
\sum_{\mathcal O\subset G_1(\Q)\cdot w}J_{\mathcal O}(f)=\tfrac12\,
\mathrm{vol}\bigl(T_w(\Q)\backslash T_w(\A)\bigr)\bigl(
J^{\mathrm{st}}_G(w,f)+J^{\mathrm{st}}_G(w,f_-)\bigr).
\]
At the real place the stable class of a lift of $w$ consists of two
conjugacy classes, exchanged by $\mathrm{Ad}(\delta_{-1})$
(Example~\ref{ex:Delta-real}); at a prime $p\equiv3\pmod4$ it consists of
two classes, of which one meets $K_p^*$; at $p\equiv1\pmod4$ it is one
class.
\end{example}

\subsection{The hyperbolic terms}\label{sec:hyperbolic-gl2}

The stable class of a hyperbolic element is its conjugacy class, so the
orbital integrals $J_G(x,f)$, $x\in\AS$ regular, are stable
distributions. The invariant distributions $I_A(x,f)$ of
Section~\ref{sec:IAgamma} are not orbital integrals, and their
stabilization is a substantial part of~\cite{LL} and of~\cite{Li21}. Here
we only determine how they transform under the untwisted action of
$G_1(k_S)$; this is the correct form of the variance formulas that the
1995 notes state for conjugation by $GL(2)$.

Let $S$, $\GS$, $\AS$ be as in Section~\ref{sec:invariant}, and
$G_1(k_S)=\prod_{v\in S_k}GL(2,k_v)$, acting on $\GS$ by $\mathrm{Ad}$.
Extend the height to $G_1$: for $g\in GL(2,k_v)$ put
\[
H_{1,v}(g)=-\log\|e_2g\|_v+\tfrac12\log|\det g|_v,
\]
so that $H_{1,v}=H_v$ on $SL(2,k_v)$ (Lemma~\ref{lem:height-ineq}(a)),
$H_{1,v}(zng)=H_{1,v}(g)$ for $z$ scalar and $n\in N$, and
$H_{1,v}(\diag(t_1,t_2)g)=\frac12\log|t_1/t_2|_v+H_{1,v}(g)$; let
$H_1=\sum_{v\in S_k}H_{1,v}$. For $\delta\in G_1(k_S)$ define
$u'_B(\tilde k,\delta)=-H_1(k\delta)$, $k=p(\tilde k)$, and $f_{B,\delta}$
by~\eqref{eq:fBy} with this weight; as in Lemma~\ref{lem:uprime}(b),
$f_{B,\delta}\in C_c^\infty(\AS)_-$. For $a\in k_S^\times$ let $\kappa_a$
be the character $\tilde a\mapsto\prod_{v\in S_k}(a_v,t_v)_v$ of $\AS$,
$p(\tilde a)=a(t)$; by Lemma~\ref{lem:Ad}(d), $\pi\circ
\mathrm{Ad}(\delta_a)=\pi\kappa_a$ for $\pi\in\Pi_u(\AS)$. We need one
more property of the normalizing factors of
Theorem~\ref{thm:normalization}:
\begin{enumerate}[label=(N\arabic*),start=5]
\item\label{N5} for every $v\in S_k$, $a\in k_v^\times$ and $\pi_v$, the
  quotient $r_v(\pi_v\kappa_a,\lambda)/r_v(\pi_v,\lambda)$ does not depend
  on $\lambda$.
\end{enumerate}
Since $\pi_v\kappa_a=\chi_{\psi,v}\cdot\chi_v(a,\cdot)_v$, the condition
holds whenever $r_v(\pi_v,\lambda)$ depends on $\pi_v=\chi_{\psi,v}\chi_v$
only through $\chi_v^2$, up to a constant factor. The explicit factors of
Lemma~\ref{lem:explicit-r} are of this form; so~\ref{N5} holds where
$\pi_v$ and $\pi_v\kappa_a$ are both covered by that lemma, which is
always the case at the real and complex places. At the other places
\ref{N5} is a condition on the choice of the factors, in addition
to~\ref{N1}--\ref{N4}. We have not proved that it can always be met
(Remark~\ref{rem:N5}). The statements below say where it is used, and
Lemma~\ref{lem:variance-gl2}(c) says what happens without it.

\begin{lemma}\label{lem:variance-gl2}
Let $f\in C_c^\infty(\GS)_-$, $a\in k_S^\times$, $\delta=\delta_a$, and
$\alpha=\mathrm{Ad}(\delta)$.
\begin{enumerate}[label=(\alph*)]
\item For $x\in\GS$ with Iwasawa decompositions $x=n\tilde a\tilde k$,
  $wx=n'\tilde a'\tilde k'$,
  \[
  v_A\bigl(\alpha^{-1}(x)\bigr)=v_A(x)-H_1(k\delta)-H_1(k'\delta).
  \]
\item For $\gamma\in\AS$ regular,
  $J_A(\gamma,f^\delta)=J_A(\alpha\gamma,f)+f_{B,\delta}(\alpha\gamma)+
  f_{B,\delta}(w\,\alpha\gamma\,w^{-1})$, and $J_G(\gamma,f^\delta)=
  J_G(\alpha\gamma,f)$.
\item For $\pi\in\Pi_u(\AS)$ and $\lambda\in i\R$, with $\pi'=\pi\circ
  \alpha^{-1}=\pi\kappa_a$: $\operatorname{tr}I(\pi_\lambda,f^\delta)=
  \operatorname{tr}I(\pi'_\lambda,f)$, and
  \[
  \varphi_A(f^\delta,\pi_\lambda)=\varphi_A(f,\pi'_\lambda)+
  \widehat{f_{B,\delta}}(\pi'_\lambda)+\widehat{f_{B,\delta}}\bigl(
  (\pi')^w_{-\lambda}\bigr)-d_a(\pi,\lambda)\,\operatorname{tr}
  I(\pi'_\lambda,f),
  \]
  where $d_a(\pi,\lambda)=\sum_{v\in S_k}\bigl(\frac{r_v'}{r_v}(\pi_v
  \kappa_a,\lambda)-\frac{r_v'}{r_v}(\pi_v,\lambda)\bigr)$. The
  condition~\ref{N5} says that $d_a=0$.
\item The products $r_v(\pi_v,\lambda)r_v(\pi_v^w,-\lambda)$ and
  $r_v(\pi_v\kappa_a,\lambda)r_v((\pi_v\kappa_a)^w,-\lambda)$ are equal up
  to a sign that does not depend on $\lambda$. So~\ref{N5} is compatible
  with~\ref{N2}.
\end{enumerate}
\end{lemma}

\begin{proof}
(a) The image of $\alpha^{-1}(x)$ in $G(\Q_S)$ is $\delta^{-1}g\delta$,
$g=p(x)$. By the properties of $H_1$, $H(\delta^{-1}g\delta)=H_1(
\delta^{-1})+H(g)+H_1(k\delta)$. Since $w\delta^{-1}w^{-1}=\diag(1,a^{-1})$,
$H(w\delta^{-1}g\delta)=H_1(\diag(1,a^{-1}))+H(wg)+H_1(k'\delta)$, and
$H_1(\diag(1,a^{-1}))=-H_1(\delta^{-1})$. Subtract from
$v_A=-H(\cdot)-H(w\,\cdot)$.

(b) $x\mapsto\alpha(x)$ preserves the Haar measure of $G(\Q_S)$ and fixes
$A(\Q_S)$ pointwise, so it preserves the measure on the quotient.
Substituting $x=\alpha^{-1}(x')$ in~\eqref{eq:JA},
\[
J_A(\gamma,f^\delta)=|D(\gamma)|^{1/2}\int f\bigl(x'^{-1}\alpha(\gamma)x'
\bigr)\,v_A\bigl(\alpha^{-1}(x')\bigr)\,dx' .
\]
Insert~(a) and argue as in the proof of
Lemma~\ref{lem:variance-local}(b), which uses only the form of the
difference of the weights.

(c) Realize $I(\pi_\lambda)$ on the space of functions $\Phi$ on $\GS$
with $\Phi(n\tilde ax)=\pi(\tilde a)e^{(\lambda+1)H_S(\tilde a)}\Phi(x)$,
by right translation, and let $(A(\lambda)\Phi)(x)=\Phi(\alpha^{-1}x)$.
Since $\alpha$ preserves $N(\Q_S)$ and $H_S$ on $\AS$, $A(\lambda)$ maps
$I(\pi_\lambda)$ to $I(\pi'_\lambda)$, and $A(\lambda)^{-1}I(\pi'_\lambda,
y)A(\lambda)=I(\pi_\lambda,\alpha^{-1}y)$; integrating against $f$,
\begin{equation}\label{eq:A-conj}
A(\lambda)^{-1}\,I(\pi'_\lambda,f)\,A(\lambda)=I(\pi_\lambda,f^\delta),
\end{equation}
which gives the first statement. In the compact picture $A(\lambda)$
depends on $\lambda$ through $e^{(\lambda+1)H_S(\alpha^{-1}\tilde k)}$, so
$\frac{d}{d\lambda}A(\lambda)=\tilde HA(\lambda)$, where $\tilde H$ is
multiplication by $\tilde k\mapsto H_S(\alpha^{-1}\tilde k)=H_1(\delta^{-1})
+H_1(k\delta)$. Next we compare intertwining operators. By
Lemma~\ref{lem:Ad}(d), $\alpha^{-1}((w,1)^{-1})=(w,1)^{-1}\tilde b$ with
$\tilde b\in\AS$ over $a(a)$, and $\alpha^{-1}(n(u))=n(a^{-1}u)$. Two
changes of variable in the integral defining $M(\pi',\lambda)$ give
\[
M(\pi',\lambda)\,A(\lambda)=c(\lambda)\,A(-\lambda)\,M(\pi,\lambda),
\qquad c(\lambda)=\pi^w(\tilde b)\,|a|_S^{-\lambda},
\]
where $A(-\lambda)$ is formed with $\pi^w$. Hence, with $R=r^{-1}M$ and
$\rho(\lambda)=r(\pi',\lambda)\,c(\lambda)^{-1}r(\pi,\lambda)^{-1}$,
\[
A(-\lambda)\,R_S(\pi,\lambda)=\rho(\lambda)\,R_S(\pi',\lambda)\,A(\lambda),
\qquad\frac{\dot\rho}{\rho}=\frac{\dot r}{r}(\pi',\lambda)-\frac{\dot r}{r}
(\pi,\lambda)+\log|a|_S ,
\]
where a dot denotes $d/d\lambda$. (The scalar $r(\pi,\lambda)$ is the product of the
$r_v(\pi_v,\lambda)$, $v\in S_k$.) Differentiate, and multiply on the left
by $(\rho R_S(\pi',\lambda))^{-1}=A(\lambda)R_S(\pi,\lambda)^{-1}
A(-\lambda)^{-1}$. Writing $R$ for $R_S(\pi,\lambda)$ and $R^\natural$ for
$R_S(\pi',\lambda)$,
\[
(R^\natural)^{-1}\dot R^\natural=A(\lambda)\,R^{-1}\dot R\,A(\lambda)^{-1}
-\frac{\dot\rho}{\rho}-\tilde H-(R^\natural)^{-1}\tilde HR^\natural .
\]
Now $\tilde H=-\frac12\log|a|_S-u'$, where $u'$ is multiplication by
$u'_B(\tilde k,\delta)=-H_1(k\delta)$, and $\dot\rho/\rho=d_a(\pi,\lambda)+
\log|a|_S$. The scalars cancel:
\[
-\frac{\dot\rho}{\rho}-\tilde H-(R^\natural)^{-1}\tilde HR^\natural
=\underbrace{-\log|a|_S+\tfrac12\log|a|_S+\tfrac12\log|a|_S}_{=\,0}
-d_a(\pi,\lambda)+u'+(R^\natural)^{-1}u'R^\natural .
\]
Take
the trace against $I(\pi'_\lambda,f)$, use~\eqref{eq:A-conj},
\ref{N1}, and Lemma~\ref{lem:uprime}(c) for the multiplication operator
$-H_1(k\delta)=u'_B(\tilde k,\delta)$:
\[
\varphi_A(f,\pi'_\lambda)=\varphi_A(f^\delta,\pi_\lambda)-
\widehat{f_{B,\delta}}(\pi'_\lambda)-\widehat{f_{B,\delta}}\bigl((\pi')^w_{
-\lambda}\bigr)+d_a(\pi,\lambda)\operatorname{tr}I(\pi'_\lambda,f).
\]

(d) Apply the relation $M(\pi',\lambda)A(\lambda)=c(\lambda)A(-\lambda)
M(\pi,\lambda)$ to $\pi$ and to $\pi^w$ at $-\lambda$; the second constant
is $\pi(\tilde b)|a|_S^{\lambda}$. So $M((\pi')^w,-\lambda)M(\pi',\lambda)$
is the conjugate by $A(\lambda)$ of $\pi(\tilde b)\pi^w(\tilde b)\cdot
M(\pi^w,-\lambda)M(\pi,\lambda)$. By~\ref{N2} the two compositions are the
scalars $r(\pi',\lambda)r((\pi')^w,-\lambda)$ and $r(\pi,\lambda)
r(\pi^w,-\lambda)$, and $\pi(\tilde b)\pi^w(\tilde b)=\pi\bigl(\tilde b\,
w\tilde bw^{-1}\bigr)$, where $\tilde b\,w\tilde bw^{-1}$ lies over $1$, is
$\pm1$. The argument is local at each $v\in S_k$.
\end{proof}

\begin{proposition}\label{prop:IA-gl2}
Assume~\ref{N5}. Let $f\in C_c^\infty(\GS)_-$ and $h\in G_1(k_S)$.
\begin{enumerate}[label=(\alph*)]
\item $I_A(x,f^h)=I_A(\mathrm{Ad}(h)x,f)$ for $x\in\AS$ regular or
  central, whenever $\mathrm{Ad}(h)x\in\AS$.
\item For $h=\delta_a$ and $x$ over $a(t)$:
  $I_A(x,f^{\delta_a})=(a,t)_S\,I_A(x,f)$, with $(a,t)_S=\prod_{v\in S_k}
  (a_v,t_v)_v$; for a lift $\tilde z$ of $-1$,
  $I_A(\tilde z,f^{\delta_a})=(a,-1)_S\,I_A(\tilde z,f)$.
\item $\operatorname{tr}I(\pi,f^{\delta_a})=\operatorname{tr}I(\pi\kappa_a,
  f)$, and $\operatorname{tr}I(\pi,f^{[h]})=\operatorname{tr}I(\pi,f)$:
  the characters of the principal series are stable. (This part does not
  use~\ref{N5}. The function $f^{[h]}$ need not be smooth; the trace is
  then defined by the integral of orbital integrals in the proof.)
\end{enumerate}
\end{proposition}

\begin{proof}
(a) For $h\in SL(2,k_S)$ this is Theorem~\ref{thm:IA}(b). Let
$h=\delta=\delta_a$. By Lemma~\ref{lem:variance-gl2}(c) and Fourier
inversion on $\AS$, substituting $\pi\mapsto\pi'$ and using
$\pi(\tilde a)=\pi'(\alpha\tilde a)$,
\[
\Phi_A(\tilde a,f^\delta)=\Phi_A(\alpha\tilde a,f)+f_{B,\delta}(\alpha
\tilde a)+f_{B,\delta}(w\,\alpha\tilde a\,w^{-1}).
\]
Subtract from Lemma~\ref{lem:variance-gl2}(b). For central $x$ use the
limit of Lemma~\ref{lem:JA-central}(b), as for
Lemma~\ref{lem:variance-local}(c). (b) follows from~(a),
Lemma~\ref{lem:Ad}(d) and the fact that $I_A(\cdot,f)$ is anti-genuine.
(c) The first statement is Lemma~\ref{lem:variance-gl2}(c). By the proof
of Theorem~\ref{thm:IA}(a), $\operatorname{tr}I(\pi,f)=\int_{\AS}\pi(\tilde
a)J_G(\tilde a,f)\,d\tilde a$, an integral of hyperbolic orbital
integrals, and $J_G(\tilde a,f^{[h]})=J_G(h\star\tilde a,f)=J_G(\tilde
a,f)$ because $h\star\tilde a$ is conjugate to $\tilde a$.
\end{proof}

\begin{remarkrev}\label{rem:N5}
\emph{The condition~\ref{N5}.} By the proof of
Lemma~\ref{lem:variance-gl2}(c), \ref{N5} holds at $v$ if and only if
\[
R_v(\pi_v\kappa_a,\lambda)=C\,|a|_v^{-\lambda}\,A(-\lambda)\,R_v(\pi_v,
\lambda)\,A(\lambda)^{-1}
\]
with a constant $C$; the factor $|a|_v^{-\lambda}$ removes the scalar
$|a|_v^{\lambda}$ contained in $A(-\lambda)A(\lambda)^{-1}$. It is natural:
the inducing data $\pi_v$ and $\pi_v\kappa_a$ have the same $\chi_v^2$,
hence the same parameter on the side of $SO(3)$
(Section~\ref{sec:stable-conj}, \emph{The dual group}), and factors defined
through that parameter, as in Lemma~\ref{lem:explicit-r}, satisfy it.
Lemma~\ref{lem:variance-gl2}(d) removes one possible conflict
with~\ref{N2}; it is a necessary compatibility and not an existence
statement. We have not proved that factors with
\ref{N1}--\ref{N5} exist at the places where no explicit factor is
available (the places of the sets $S_{\mathrm{bad}}(\mu)$ of
Section~\ref{sec:normalized}), nor checked
whether the factors of~\cite{LiII} have property~\ref{N5}. One cannot
simply transport $r_v$ along the orbits of $k_v^\times/k_v^{\times2}$ on
the genuine characters: when $\pi_v^w=\pi_v\kappa_{a_0}$ with $a_0$ not a
square, that is, when $\chi_v^2=(a_0,\cdot)_v\ne1$, the relation~\ref{N2}
between $r_v(\pi_v,\cdot)$ and $r_v(\pi_v^w,\cdot)$ constrains the factor
of $\pi_v$ itself. So in this section \ref{N5} is a hypothesis at the
places where the factors are not explicit. What would remove it is the
following statement, which we formulate as a problem: \emph{for every
local field $F$ of characteristic $0$, including the dyadic ones, and
every $\psi$, there are functions $F_v(\chi^2,\lambda)$ and constants
$C_v(\chi)$ such that the factors $r_v(\chi_{\psi}\chi,\lambda)=C_v(\chi)
F_v(\chi^2,\lambda)$ satisfy \ref{N1}--\ref{N4}.} The literature we know
does not contain it. The local coefficients of~\cite{Sz09} are computed in
all residue characteristics, but we found no statement there that gives
\ref{N1}--\ref{N4} for factors of this form. The Plancherel measures
of~\cite{GS16} are computed under the assumption that the residue
characteristic is prime to the degree of the cover, which excludes the
dyadic places for the double cover. The paper~\cite{Sz16} drops that
assumption when $4\nmid n$; we have not worked out what it gives here.
Without~\ref{N5}, every formula
below that is stated under~\ref{N5} holds with the additional term
coming from $d_a$ in Lemma~\ref{lem:variance-gl2}(c). The distributions
$I_A(x,f)$ depend on the choice of the factors in any case
(Remark~\ref{rem:bad-places}).
\end{remarkrev}

So $I_A(x,\cdot)$ transforms under $G_1(k_S)$ by the character
$h\mapsto(\det h,t)_S=\eta_x(\det h)$, as the orbital integral
$J_G(x,\cdot)$ does. This does \emph{not} say that $I_A(x,\cdot)$ is
stable. Stability is invariance under $f\mapsto f^{[h]}$, where the sign
$\eta_y(\det h)$ varies with the point $y$ of the support of $f$, and
$I_A(x,f)$ depends on $f$ away from the class of $x$. At a real place, for instance, $\Phi_A(x,f)$ is built from the operators
$R_v^{-1}R_v'$, whose poles are the points $\lambda=a_i,b_i$ at which
$I(\pi_\lambda)$ has a single discrete series representation as a
subquotient (Example~\ref{ex:real-R}), while a stable character involves
two discrete series representations with different central
characters~\cite{Ad}. The stabilization of these terms in~\cite{Li21}
expresses them through stable distributions on the endoscopic groups for
both endoscopic data; it requires the weighted fundamental
lemma~\cite{Li12} and is not attempted here.

\subsection{The central terms and the unipotent coefficients}
\label{sec:unipotent}

The class $\mathcal O_z$ of $z=\pm1$ contributes $J_{\mathrm{unip}}(f_z)$,
which Proposition~\ref{prop:central} expresses through the zeta integral
$Z(f,s)$. Its second form in~\eqref{eq:Zf} is a sum over the rational
unipotent classes $\xi\in k^\times/k^{\times2}$; grouped according to the
image of $\xi$ in $k_S^\times/k_S^{\times2}$ it diverges at $s=1$, class by
class; this is the difficulty noted in
Remark~\ref{rem:sec4-status}. The way out, as in~\cite{LL}, is to expand
in the characters of the compact group $\A_k^\times/k^\times\A_k^{\times2}$
instead. This computes the coefficients $a^G(S,u_\xi)$ of
Theorem~\ref{thm:fine}.

Let $S$ be as in Theorem~\ref{thm:tf}: a finite set of places of $\Q$
containing $\infty$, $2$ and the places below which $\psi$ is ramified.
In this subsection we assume in addition that the ring $\OO_S$ of
elements of $k$ that are integral outside $S_k$ is a principal ideal
domain, which holds once $S$ is large enough (and always for $k=\Q$).
Let $\mathfrak X_S$ be the group of quadratic characters $\kappa$ of
$k^\times\backslash\A_k^\times$ that are unramified outside $S_k$,
$\kappa_S$ the restriction of $\kappa$ to $k_S^\times$, and
$L^S(s,\kappa)=\prod_{v\notin S_k}(1-\kappa_v(\varpi_v)q_v^{-s})^{-1}$. For
$\kappa=1$ this is $\zeta_k^S(s)$; write
\[
\zeta_k^S(s)=\frac{\rho_S}{s-1}+c_S+O(s-1).
\]
By Lemma~\ref{lem:w-action}(d) and its proof, $\mathfrak X_S=\{\kappa_b
\mid b\in\OO_S^\times/\OO_S^{\times2}\}$ with $\kappa_b(t)=\prod_v(b,t_v)_v$;
it has $2^{|S_k|}$ elements, by the unit theorem. For $f_S=\bigotimes_{v
\in S_k}f_v$, $z\in\{\pm1\}$ and $\Phi_v$ as in
Lemma~\ref{lem:JA-central}, put
\[
Z_v(\Phi_v,\kappa_v,s)=\int_{k_v^\times}\Phi_v(y)\kappa_v(y)|y|_v^s\,d^\times y,
\qquad Z_S(\Phi_S,\kappa_S,s)=\prod_{v\in S_k}Z_v(\Phi_v,\kappa_v,s),
\]
Tate's local zeta integrals~\cite{Ta}, holomorphic for $\mathrm{Re}\,s>0$;
by linearity $Z_S$ is defined for all $f_S\in\mathcal H(\GS)$. For
$\xi\in k_S^\times$ let
\begin{equation}\label{eq:JG-unip}
J_G(zu_\xi,f_S)=|\xi|_S\int_{\{\pm1\}N(\Q_S)\backslash G(\Q_S)}
f_S\bigl(x^{-1}z\,n(\xi)\,x\bigr)\,dx
=\prod_{v\in S_k}\frac{|\xi|_v}2\int_{k_v^\times}\Phi_v(\xi t^2)|t|_v^2\,
d^\times t,
\end{equation}
the second expression by Lemma~\ref{lem:JA-central}(c). The factor
$|\xi|_S$ makes $J_G(zu_\xi,f_S)$ depend only on the class of $\xi$ in
$k_S^\times/k_S^{\times2}$; it is the orbital integral for the measure on
the centralizer of $n(\xi)$ transported from that of $n(1)$.

\begin{proposition}\label{prop:unipotent}
Let $S$ be as above (in particular the ring $\OO_S$ of $S_k$-integers is
principal), $f_S\in\mathcal H(\GS)$, $f=f_S\otimes e^S$ and
$z\in\{\pm1\}$.
\begin{enumerate}[label=(\alph*)]
\item For $\mathrm{Re}\,s>1$,
  \[
  Z(f_z,s)=\frac12\sum_{\kappa\in\mathfrak X_S}L^S(s,\kappa)\,
  Z_S(\Phi_S,\kappa_S,s).
  \]
\item For every character $\kappa_S$ of $k_S^\times/k_S^{\times2}$,
  \[
  Z_S(\Phi_S,\kappa_S,1)=\sum_{\xi\in k_S^\times/k_S^{\times2}}\kappa_S(\xi)\,
  J_G(zu_\xi,f_S)=:J^\kappa_{\mathrm{unip}}(z,f_S),
  \]
  and $\frac{d}{ds}Z_S(\Phi_S,1,s)\big|_{s=1}=\rho_S^{-1}c_k\,J_A(z,f_S)$.
\item Consequently
  \begin{multline*}
  J_{\mathcal O_z}(f)=\mathrm{vol}\bigl(G(\Q)\backslash\GA\bigr)f_S(z)
  +\frac{c_k}2\,J_A(z,f_S)\\
  +\frac{c_S}2\,J^1_{\mathrm{unip}}(z,f_S)
  +\frac12\sum_{\kappa\in\mathfrak X_S,\ \kappa\ne1}L^S(1,\kappa)\,
  J^\kappa_{\mathrm{unip}}(z,f_S).
  \end{multline*}
  This is Theorem~\ref{thm:fine}(a), with
  \[
  a^G(S,u_\xi)=\frac12\Bigl(c_S+\sum_{\kappa\in\mathfrak X_S,\ \kappa\ne1}
  \kappa_S(\xi)\,L^S(1,\kappa)\Bigr),\qquad\xi\in k_S^\times/k_S^{\times2}.
  \]
\item For $a\in k_S^\times$,
  \[
  J_G(zu_\xi,f_S^{\delta_a})=(a,z)_S\,J_G(zu_{a\xi},f_S),\qquad
  J^\kappa_{\mathrm{unip}}(z,f_S^{\delta_a})=\kappa_S(a)\,(a,z)_S\,
  J^\kappa_{\mathrm{unip}}(z,f_S).
  \]
\end{enumerate}
\end{proposition}

\begin{proof}
We may replace $f$ by $f_z$ and take $z=1$; then $\Phi_f=\Phi_S\otimes
\bigotimes_{v\notin S_k}\mathbf 1_{\OO_v}$.

\emph{A local identity.} For $h$ integrable on $k_v^\times$ and
$\xi\in k_v^\times$,
\begin{equation}\label{eq:squares}
\int_{k_v^\times}h(\xi t^2)\,d^\times t=\frac12\sum_{\kappa_v}\kappa_v(\xi)
\int_{k_v^\times}h(y)\kappa_v(y)\,d^\times y,
\end{equation}
the sum over the characters of $k_v^\times/k_v^{\times2}$. Indeed
$t\mapsto t^2$ is two-to-one onto the open subgroup $k_v^{\times2}$ and
multiplies $d^\times t$ by $|2|_v$, so the left side is
$2|2|_v^{-1}\int_{\xi k_v^{\times2}}h\,d^\times y$; and the characteristic
function of $\xi k_v^{\times2}$ is $[k_v^\times:k_v^{\times2}]^{-1}
\sum_{\kappa_v}\kappa_v(\xi)\kappa_v$. The index is $4|2|_v^{-1}$ at every
place ($4|2|_v^{-1}$ for $v$ finite, $2$ for $v$ real, $1$ for $v$
complex, with the normalized absolute values).

(a) Let $\Gamma=\OO_S^\times$ and $Y_1=k_S^\times\times\prod_{v\notin S_k}
\OO_v^\times$. Since $\OO_S$ is principal, $\A_k^\times=k^\times Y_1$ and
$k^\times\cap Y_1=\Gamma$, so $\A_k^\times$ is the disjoint union of the
$cY_1$, $c\in k^\times/\Gamma$. For $\eta\in k^\times$ let
$J(\eta)=\int_{Y_1}\Phi_f(\eta t^2)|t|^{2s}d^\times t$. Then the integral
in the second form of~\eqref{eq:Zf} is $\sum_{c\in k^\times/\Gamma}J(\xi
c^2)$, and $(\xi,c)\mapsto\xi c^2$ is a bijection from
$(k^\times/k^{\times2})\times(k^\times/\Gamma)$ onto $k^\times/\Gamma^2$
(because $-1\in\Gamma$). So $Z(f,s)=\frac12\sum_{\eta\in k^\times/\Gamma^2}
J(\eta)$. Now $J(\eta)=0$ unless $\eta\in\OO_S$, and
by~\eqref{eq:squares} at each $v\in S_k$, applied to $h=\Phi_v|\cdot|_v^s$,
\[
J(\eta)=|\eta|_S^{-s}\,2^{-|S_k|}\sum_{\kappa_S}\kappa_S(\eta)\,
Z_S(\Phi_S,\kappa_S,s)\qquad(\eta\in\OO_S\smallsetminus\{0\}).
\]
Summing over $\eta$ modulo $\Gamma^2$: the sum of $\kappa_S$ over
$\Gamma/\Gamma^2$ is $2^{|S_k|}$ if $\kappa_S$ is trivial on $\Gamma$ and
$0$ otherwise. The characters $\kappa_S$ trivial on $\Gamma$ are the
restrictions of the $\kappa\in\mathfrak X_S$, and for such $\kappa$ and
$\eta\in\OO_S$, $\kappa_S(\eta)=\prod_{v\notin S_k}\kappa_v(\varpi_v)^{
\ord_v\eta}$ and $|\eta|_S^{-1}=\prod_{v\notin S_k}q_v^{\ord_v\eta}$, by the
product formula. The sum over $\eta\in(\OO_S\smallsetminus\{0\})/\Gamma$
runs over the nonzero ideals of $\OO_S$ and gives $L^S(s,\kappa)$.

(b) By~\eqref{eq:squares} with $h=\Phi_v|\cdot|_v$ and~\eqref{eq:JG-unip},
$J_{G,v}(u_\xi,f_v)=\frac14\sum_{\kappa_v}\kappa_v(\xi)Z_v(\Phi_v,\kappa_v,
1)$; invert on $k_v^\times/k_v^{\times2}$ and multiply over $v\in S_k$. The
constant is $\prod_{v\in S_k}4\,[k_v^\times:k_v^{\times2}]^{-1}=\prod_{v\in
S_k}|2|_v=1$, because $S_k$ contains all places at which $|2|_v\ne1$. For
the derivative, let $d^\times t=c'_v\,dt/|t|_v$ at $v\in S_k$, $dt$ the
self-dual measure, and $C'=\prod_{v\in S_k}c'_v$. Then
$Z_S(\Phi_S,1,1)=C'\prod_v\hat\Phi_v(0)$ and
\[
\frac{d}{ds}Z_S(\Phi_S,1,s)\Big|_{s=1}=C'\sum_{v\in S_k}\Bigl(\int_{k_v}
\Phi_v(y)\log|y|_v\,dy\Bigr)\prod_{w\ne v}\hat\Phi_w(0)=C'\,J_A(1,f_S),
\]
by Lemma~\ref{lem:JA-central}(b), since $\sum_{v\in S_k}\log|2|_v=0$.
Comparing the residue of~(a) at $s=1$, which is $\frac12\rho_SC'
\prod_v\hat\Phi_v(0)$, with the residue $\frac12c_k\hat\Phi_f(0)$ of
Proposition~\ref{prop:central} gives $\rho_SC'=c_k$.

(c) By~(a), the finite part of $Z(f,s)$ at $s=1$ is $\frac12\bigl(c_S
Z_S(\Phi_S,1,1)+\rho_S\frac{d}{ds}Z_S(\Phi_S,1,s)|_{s=1}\bigr)+\frac12
\sum_{\kappa\ne1}L^S(1,\kappa)Z_S(\Phi_S,\kappa_S,1)$; for $\kappa\ne1$,
$L^S(s,\kappa)$ is holomorphic at $s=1$. Insert~(b) and
Proposition~\ref{prop:central}. The coefficient of $J_G(zu_\xi,f_S)$ is
read off from the definition of $J^\kappa_{\mathrm{unip}}$.

(d) $x\mapsto\delta_ax\delta_a^{-1}$ preserves the Haar measure of
$G(\Q_S)$ and multiplies that of $N(\Q_S)$ by $|a|_S$; so
$\int h(\delta_ax\delta_a^{-1})\,dx=|a|_S\int h(x)\,dx$ for the invariant
measure on the quotient by $\{\pm1\}N(\Q_S)$. With Lemma~\ref{lem:Ad}(d),
this gives $\int f_S^{\delta_a}(
x^{-1}zn(\xi)x)\,dx=(a,z)_S\,|a|_S\int f_S(x^{-1}zn(a\xi)x)\,dx$. Multiply
by $|\xi|_S$, and sum against $\kappa_S(\xi)$.
\end{proof}

Part~(c) proves the fine expansion at the central classes directly from
the coarse one, for the sets $S$ considered, and so removes the appeal
to~\cite{A12,LiI} in Theorem~\ref{thm:fine}(a). The terms with
$\kappa\ne1$ are the analog of the terms of~\cite[\S5]{LL} that involve
$L(1,\kappa)$; in the 1995 notes they were written as a sum over the
rational unipotent classes, which diverges
(Appendix~\ref{sec:appendix-1995}).

\begin{example}\label{ex:unip-coeff}
Let $k=\Q$ and $S=\{\infty,2\}$. Then $\OO_S=\Z[\frac12]$,
$\mathfrak X_S=\{\kappa_b\mid b=\pm1,\pm2\}$, and
$\zeta^S(s)=(1-2^{-s})\zeta(s)$, so $\rho_S=\frac12$ and
$c_S=\frac12(\gamma_E+\log2)$, with Euler's constant $\gamma_E$. The
characters $\kappa_{-1}$, $\kappa_2$, $\kappa_{-2}$ are the Dirichlet
characters of conductor $4$, $8$, $8$, and
\[
L^S(1,\kappa_{-1})=\frac\pi4,\qquad L^S(1,\kappa_2)=\frac{\log(1+\sqrt2)}
{\sqrt2},\qquad L^S(1,\kappa_{-2})=\frac\pi{2\sqrt2}.
\]
For the $16$ classes $\xi\in\Q_S^\times/\Q_S^{\times2}$,
\[
a^G(S,u_\xi)=\frac12\Bigl(\frac{\gamma_E+\log2}2+\sum_{b\in\{-1,2,-2\}}
(b,\xi_\infty)_\infty(b,\xi_2)_2\,L^S(1,\kappa_b)\Bigr);
\]
for instance $a^G(S,u_\xi)\approx1.57726$ for $\xi=1$, and
$a^G(S,u_\xi)\approx-0.31886$ for the class with $\xi_\infty<0$ and $\xi_2=1$.
The coefficients are not all positive, and they depend on $\xi$ only through
the pairings with $b\in\{\pm1,\pm2\}$, the same parameters as those of the
ramified classes $\chi_{\psi_b}$ in Example~\ref{ex:ramified-Q}. In
particular $a^G(S,u_{\xi\epsilon})=a^G(S,u_\xi)$ for $\epsilon\in\OO_S^\times$,
since the $\kappa\in\mathfrak X_S$ are trivial on $\OO_S^\times$: the
classes of the rational elements $\pm1$, $\pm2$ all have the coefficient of
$\xi=1$, and the $16$ classes carry four distinct values (the invariance
holds for every $k$ and $S$ as above). With $\epsilon_b=(b,\xi_\infty)_\infty
(b,\xi_2)_2$ they are
\begin{center}
\begin{tabular}{ccc}
$(\epsilon_{-1},\epsilon_2,\epsilon_{-2})$ & a class $\xi=(\xi_\infty,\xi_2)$ &
$a^G(S,u_\xi)$\\\hline
$(+,+,+)$ & $(1,1)$ & $\phantom{-}1.57726$\\
$(-,-,+)$ & $(1,3)$ & $\phantom{-}0.16864$\\
$(+,-,-)$ & $(1,5)$ & $-0.15668$\\
$(-,+,-)$ & $(-1,1)$ & $-0.31886$
\end{tabular}
\end{center}
Each value is taken on four classes, and the four values add up to
$2c_S=\gamma_E+\log2$. The
script checks~(a) numerically at $s=3$ for a test function with
$\Phi_\infty$ not even and $\Phi_2$ the characteristic function of
$6+16\Z_2$.
\end{example}

By part~(d), $J^\kappa_{\mathrm{unip}}(1,\cdot)$ transforms under
$G_1(k_S)$ by $\kappa_S(\det h)$. Near the identity $\eta_x=1$ (because
$D^+$ is close to $4$), so there the twisted and the untwisted actions
agree: $J^1_{\mathrm{unip}}(1,\cdot)$, the sum over all unipotent classes,
is the stable unipotent distribution, and the
$J^\kappa_{\mathrm{unip}}(1,\cdot)$ with $\kappa\ne1$ are the unstable
ones. For $SL(2)$ the latter are the limits, as $\gamma\to1$ in the torus
$T_E$ with $\kappa_E=\kappa$, of the $\kappa$-orbital integrals, up to the
factor $L(1,\kappa)$~\cite[\S\S2,~5]{LL}. On the cover the same should
hold near the identity, where the cover is canonically split on the
topologically unipotent set, and by Lemma~\ref{lem:kappa-orbital}(c) the
$\kappa$-orbital integrals near $1$ are the stable orbital integrals of
$f_-$ near $-1$. So the terms $L^S(1,\kappa)J^\kappa_{\mathrm{unip}}(1,
f_S)$, $\kappa\ne1$, should be counted with $S(f_-)$, and the
corresponding terms at $z=-1$ with $S(f)$. We have not proved these limit
formulas for the cover (Remark~\ref{rem:sec7-status}).

\subsection{The spectral side}\label{sec:spectral-gl2}

On the spectral side the group $G_1(k)$ acts on the space of automorphic
forms, and the pieces of the spectral decomposition are permuted. We
record this, and then compute the one family of discrete terms whose
stable structure is visible without endoscopy: the theta functions.

For $h\in G_1(k)$ and a function $\varphi$ on $G(\Q)\backslash\GA$ put
$(U_h\varphi)(x)=\varphi(\mathrm{Ad}(h)^{-1}x)$; this is well defined by
Lemma~\ref{lem:Ad}(f). For $a\in k^\times$ let $\kappa_a$ also denote the
character $\tilde a\mapsto\prod_v(a,t_v)_v$ of $\overline{A(\A)}$.

\begin{lemma}\label{lem:spectrum-gl2}
Let $h\in G_1(k)$ and $a=\det h$.
\begin{enumerate}[label=(\alph*)]
\item $U_h$ is a unitary operator on $L^2_{\mathrm{gen}}$ with
  $U_hR(y)U_h^{-1}=R(\mathrm{Ad}(h)y)$ and $U_hR(f)U_h^{-1}=R(f^{h^{-1}})$.
  It preserves $L^2_{\mathrm{cusp}}$, $L^2_{\mathrm{res}}$ and
  $L^2_{\mathrm{cont}}$.
\item $U_h$ maps $L^2_\chi$ onto $L^2_{\chi^h}$, where $(G,\pi)^h=
  (G,\pi\circ\mathrm{Ad}(h)^{-1})$ and $(A,\mu)^h=(A,\mu\kappa_a)$. In
  particular $m(\pi\circ\mathrm{Ad}(h))=m(\pi)$, $\chi_{\psi_b}\kappa_a=
  \chi_{\psi_{ab}}$, and $U_hV_b=V_{ab}$.
\item Let $f=f_S\otimes e^S$ and $h_v\in Z(k_v)GL(2,\OO_v)$ for all
  $v\notin S_k$. Then $\operatorname{tr}(R(f^h)|_{L^2_{\chi,\mathrm{disc}}})=
  \operatorname{tr}(R(f)|_{L^2_{\chi^h,\mathrm{disc}}})$ for every $\chi$, and
  if~\ref{N5} holds, $I_\chi(f^h)=I_{\chi^h}(f)$. Hence
  $\sum_\chi I_\chi(f^h)=\sum_\chi I_\chi(f)$.
\end{enumerate}
\end{lemma}

\begin{proof}
(a) $\mathrm{Ad}(h)$ preserves $G(\Q)$ and the Tamagawa measure, which
gives unitarity and the two formulas. It preserves $N(\A)$ and its Haar
measure up to the factor $|a|_{\A_k}=1$, so $(U_h\varphi)_B=U_h(\varphi_B)$
for the constant term along $B$; hence cusp forms go to cusp forms, and
the orthogonal complement, the discrete and the continuous parts of it,
are preserved. (b) By Lemma~\ref{lem:Ad}(d), $\mathrm{Ad}(h)^{-1}$ acts
on $\overline{A(\A)}$ by the character $\kappa_a$ (for $h=\delta_a$, and
conjugation by $SL(2,k)$ does not change the class of a cuspidal datum).
So $U_h$ maps $\mathcal H_B(\mu)$ to $\mathcal H_B(\mu\kappa_a)$ and
Eisenstein series, their wave packets and their residues to the
corresponding objects for $\mu\kappa_a$. The formula for
$\chi_{\psi_b}\kappa_a$ is Lemma~\ref{lem:weil}(d). (c) By
Lemma~\ref{lem:Ad}(e), $f^h=f_S^{h_S}\otimes e^S$. The first statement
follows from~(a) and~(b): $\operatorname{tr}(R(f^h)|_{L})=
\operatorname{tr}(R(f)|_{U_hL})$ for an invariant subspace $L$. This gives
$I_\chi(f^h)=I_{\chi^h}(f)$ for $\chi\in\mathcal X_G$ and for the residual
terms. For $\mu\in\Pi_{\mathrm{aut}}(X_A)$, $\operatorname{tr}I(\mu_{it},
f^h)=\operatorname{tr}I((\mu\kappa_a)_{it},f)$ by
Lemma~\ref{lem:variance-gl2}(c), and $\frac{m'}m(\mu\kappa_a,\lambda)=
\frac{m'}m(\mu,\lambda)$: in~\eqref{eq:global-m} the partial $L$-functions
depend only on $\chi^2$, and the local factors change by constants,
by~\ref{N5}. So the terms of Theorem~\ref{thm:itf}(e),(f) for $\chi$ and
$f^h$ are those for $\chi^h$ and $f$.
\end{proof}

The geometric counterpart of~(c) is Theorem~\ref{thm:stabler}(c) below.

Now fix $b\in k^\times$ and put $\psi'=\psi_{-b}$, as in
Theorem~\ref{thm:residual}(c). For a pattern of signs
$\varepsilon=(\varepsilon_v)_v$, with $\varepsilon_v=+$ for almost all
$v$, let $r^\varepsilon_{\psi'}=\bigotimes'_vr^{\varepsilon_v}_{\psi'_v}$, an
irreducible genuine unitary representation of $\GA$ on the functions in
$\mathcal S(\A_k)$ of parity $\varepsilon_v$ at each $v$. Let
\[
\Theta_b=\bigl\{\theta^{\psi'}_\phi\ \big|\ \phi\in\mathcal S(\A_k)\bigr\}^{-}
\subset L^2_{\mathrm{gen}},
\]
the closure of the space of all theta functions~\eqref{eq:theta}, even or
not; and for $f=\bigotimes f_v$ let
\[
\Theta^{\mathrm{st}}_b(f)=\prod_v\bigl(\operatorname{tr}r^+_{\psi'_v}(f_v)-
\operatorname{tr}r^-_{\psi'_v}(f_v)\bigr)=\prod_v\int f_v(x)\,
\frac{\Delta_{\psi'_v}(x)}{|D^+(x)|_v^{1/2}}\,dx,
\]
the second expression by Proposition~\ref{prop:weil-char}. Almost all
factors are $1$. Both $\Theta_b$ and $\Theta^{\mathrm{st}}_b$ depend only
on the class of $b$ in $k^\times/k^{\times2}$: the substitution
$\phi(x)\mapsto\phi(cx)$, $c\in k^\times$, carries $\omega_{\psi'}$ to
$\omega_{\psi'_{c^{-2}}}$, preserves parities, and does not change
$\sum_{\xi\in k}$.

\begin{proposition}\label{prop:theta-packet}
\begin{enumerate}[label=(\alph*)]
\item $\phi\mapsto\theta^{\psi'}_\phi$ is zero on $r^\varepsilon_{\psi'}$ if
  $\prod_v\varepsilon_v=-1$ and injective if $\prod_v\varepsilon_v=+1$. So
  \[
  \Theta_b\cong\bigoplus_{\prod_v\varepsilon_v=+1}r^\varepsilon_{\psi'},
  \]
  a subspace of $L^2_{\mathrm{disc}}$; the summand with all
  $\varepsilon_v=+$ is $V_b$, and the others are cuspidal.
\item For $f=\bigotimes_vf_v\in\mathcal H(\GA)$,
  \[
  \operatorname{tr}\bigl(R(f)\big|_{\Theta_b}\bigr)=\tfrac12\,
  \Theta^{\mathrm{st}}_b(f)+\tfrac12\,\Theta^{\mathrm{st}}_b(f_-).
  \]
\item $\Theta^{\mathrm{st}}_b$ is a stable distribution:
  $\Theta^{\mathrm{st}}_b(f^{[h]})=\Theta^{\mathrm{st}}_b(f)$ for
  $h\in G_1(\A_k)$. Moreover $U_h\Theta_b=\Theta_{ab}$ and
  $\Theta^{\mathrm{st}}_b(f^h)=\Theta^{\mathrm{st}}_{ab}(f)$ for
  $h\in G_1(k)$, $a=\det h$.
\end{enumerate}
\end{proposition}

\begin{proof}
(a) $\theta_{\omega(y)\phi}(x)=\theta_\phi(xy)$, so $\phi\mapsto
\theta_\phi$ intertwines $\omega_{\psi'}$ with $R$. The parity operator
$P=\bigotimes_vP_v$ satisfies $\theta_{P\phi}=\theta_\phi$ (replace $\xi$
by $-\xi$ in~\eqref{eq:theta}; $P$ commutes with $\omega$). If
$\prod\varepsilon_v=-1$, then $P\phi=-\phi$ on $r^\varepsilon$ and
$\theta_\phi=0$. If $\prod\varepsilon_v=+1$, let $S_0$ be a finite set of places
containing the archimedean ones and those with $\varepsilon_v=-$, and
$\xi_0=1$. Put $\phi_v=\mathbf 1_{\OO_v}$ for $v\notin S_0$, and
$\phi_v=h_v+\varepsilon_vP_vh_v$ for $v\in S_0$, where $h_v$ is supported
in a neighborhood $U_v$ of $\xi_0$ with $U_v\cap(-U_v)=\emptyset$ and
$h_v(\xi_0)=1$. Then $\phi=\bigotimes\phi_v\in r^\varepsilon$. Since $k$ is
discrete in $\A_k$, for $U_v$ small enough the only $\xi\in k$ that are
integral outside $S_0$ and lie in $U_v\cup(-U_v)$ for every $v\in S_0$ are
$\pm\xi_0$, and $\theta_\phi(1)=\phi(\xi_0)+\phi(-\xi_0)=1+\prod_{v\in S_0}
\varepsilon_v=2$. So the map is not zero on the irreducible
$r^\varepsilon$, hence injective. The constant term of $\theta_\phi$ is
$x\mapsto(\omega(x)\phi)(0)$ (Section~\ref{sec:eisenstein}), which
vanishes when some $\phi_v$ is odd; so those $\theta_\phi$ are cuspidal,
and the even ones span $V_b$ by Theorem~\ref{thm:residual}(c). That theta
functions are square integrable, and that $\phi\mapsto\theta_\phi$ is
continuous for the $L^2$ norms on each $r^\varepsilon$, is
standard~\cite{We64,Ge2,GPS80}: on a Siegel set $\theta_\phi(n\tilde a
\tilde k)=O(e^{H(\tilde a)/2})$. The representations $r^\varepsilon$ are
pairwise inequivalent, because $r^+_{\psi'_v}$ and $r^-_{\psi'_v}$ have
different central characters; so the sum is direct and orthogonal.

(b) By~(a), $\operatorname{tr}(R(f)|_{\Theta_b})=\sum_{\prod\varepsilon_v=1}
\prod_v\operatorname{tr}r^{\varepsilon_v}_{\psi'_v}(f_v)$, a finite sum
because $\operatorname{tr}r^-_{\psi'_v}(e_v)=0$ at the unramified places.
With $t^\pm_v=\operatorname{tr}r^\pm_{\psi'_v}(f_v)$ this is
$\frac12\prod_v(t^+_v-t^-_v)+\frac12\prod_v(t^+_v+t^-_v)$. The first
product is $\Theta^{\mathrm{st}}_b(f)$, and the second is
$\operatorname{tr}\omega_{\psi'}(f)$. By Lemma~\ref{lem:weil-rep},
$\omega_{\psi'_v}(-1,1)=\gamma'_v(1)^2P_v$, and $\prod_v\gamma'_v(1)=1$
(Lemma~\ref{lem:weil}(b)); so $\omega_{\psi'}(s(-1))=P$. Therefore
$\operatorname{tr}\omega_{\psi'}(f)=\operatorname{tr}\bigl(\omega_{\psi'}(
s(-1))\,\omega_{\psi'}(f)\,P\bigr)=\operatorname{tr}\bigl(\omega_{\psi'}
(f_-)P\bigr)=\Theta^{\mathrm{st}}_b(f_-)$, since $\omega(z)\omega(f)=
\omega(f_{z^{-1}})$.

(c) Stability is Proposition~\ref{prop:weil-char} at each place. By
Lemma~\ref{lem:weil-kernel}, $\omega_{\psi'}\circ\mathrm{Ad}(\delta_a)=
\omega_{\psi'_a}$ and $\psi'_a=\psi_{-ab}$; for $a\in k^\times$ the
identity map of $\mathcal S(\A_k)$ therefore carries $\theta^{\psi'}_\phi
\circ\mathrm{Ad}(\delta_a)$ to $\theta^{\psi'_a}_\phi$, which gives
both statements (replace $a$ by $a^{-1}$, which is in the same square
class).
\end{proof}

Part~(b) has the same two-term form as Theorem~\ref{thm:elliptic-stable}: a
stable distribution evaluated at $f$ and at $f_-$, each with coefficient
$\frac12$. The multiplicity of $r^\varepsilon_{\psi'}$ in $\Theta_b$ is
$\frac12(1+\prod_v\varepsilon_v)$. Waldspurger~\cite{Wa80,Wa91} proved
that this is its multiplicity in $L^2_{\mathrm{disc}}$, and that the rest
of the discrete spectrum has a similar structure: it is a
multiplicity-free sum over the cuspidal automorphic representations
$\sigma$ of $PGL(2,\A_k)$, each contributing the representations
$\bigotimes_v\pi_v^{\varepsilon_v}$ of a packet with $\prod_v\varepsilon_v=
\varepsilon(\frac12,\sigma)$, where the local packet is $\{\pi_v^+\}$ or
$\{\pi_v^+,\pi_v^-\}$. Here the packets are those of the theta
correspondence for the additive character $\psi'$: $\pi_v^+$ is the lift
of $\sigma_v$ from $PGL(2,k_v)=SO(2,1)$, and $\pi_v^-$, when it exists, is
the lift of the Jacquet--Langlands transfer of $\sigma_v$ from the
anisotropic $SO(3)$; the labels depend on $\psi'$. See~\cite{GS12} for the
local correspondence and its epsilon dichotomy, \cite{GI} for the double
cover of $Sp(2n)$, and~\cite{Li24}. In these terms $\Theta_b$
is the packet attached to the character $\kappa_b$ of $PGL(2)$. One
difference should be noted. For the theta functions the stable
combination is the difference $\operatorname{tr}r^+-\operatorname{tr}r^-$
(Proposition~\ref{prop:weil-char}). For the packets of cuspidal $\sigma$
one expects it to be the sum $\operatorname{tr}\pi^+_v+\operatorname{tr}
\pi^-_v$, with the difference unstable. We do not know a reference that
states this for all packets in terms of the action $\star$ used here, and
we do not prove it; it is the local character identity of the endoscopic
transfer~\cite{Li19}. Two cases can be checked directly. At a finite
place, the packet of the Steinberg representation is
$\{\mathrm{sp},r^-_{\psi'}\}$~\cite[\S5]{GS12},
where $\mathrm{sp}$ is the other constituent of the principal series
that has $r^+_{\psi'}$ as a quotient (it is the lift of the trivial
representation of the anisotropic $SO(3)$, and $r^-_{\psi'}$ is the lift
of the Steinberg representation); the sum of their characters is the
character of that principal series minus $\Theta^+_{\psi'}-\Theta^-_{\psi'}$,
a difference of two stable characters
(Propositions~\ref{prop:IA-gl2}(c) and~\ref{prop:weil-char}). At a real
place a packet consists of the discrete series representations with
weights $m,m+2,\dots$ and $-m,-m-2,\dots$, $m\in\frac12+\Z$, $m>1$
(see~\cite[\S15]{Ad} for the stable discrete series characters of
$Mp(2,\R)$ and the central characters of the two members); on
$\tilde k_\theta$ the sum of their characters is $-\sin((m-1)\theta)/
\sin\theta$, which is even in $\theta$, and the difference is
$i\cos((m-1)\theta)/\sin\theta$, which is odd, while $\tilde k_\theta$ and
$\tilde k_{-\theta}$ are stably conjugate (Example~\ref{ex:Delta-real}).
In both kinds of packet the two members have opposite central characters,
so the unstable combination for $f$ is, up to a constant, the stable one
for $f_-$; this is how the pattern $S(f)+S(f_-)$ arises on the spectral
side. We do not derive any of this from the trace formula. The
multiplicity statements are quoted from~\cite{Wa80,Wa91} as recounted
in~\cite{GS12}; we have not checked them against the original papers.

By Proposition~\ref{prop:IA-gl2}(c) the terms of
Corollary~\ref{cor:itf-final} that come from the continuous spectrum,
$-\frac14\operatorname{tr}I(\chi_{\psi_b,0},f)$ and $-\frac1{4\pi}\int
\frac{m'}m(\mu,it)\operatorname{tr}I(\mu_{it},f)\,dt$, are stable
distributions. They are unchanged when $f$ is replaced by $f_-$, because
$I(\mu_\lambda,s(-1))=\mu(s(-1))=1$ for $\mu\in\Pi_{\mathrm{aut}}(X_A)$.

\subsection{The stable(r) form}\label{sec:stable-form}

We collect the results. For $f=f_S\otimes e^S$ as in
Theorem~\ref{thm:itf}, with $S$ as in Section~\ref{sec:unipotent}, write
$I(f)=\sum_{\mathcal O}I_{\mathcal O}(f)=\sum_\chi I_\chi(f)$.

Not everything in the theorem depends on~\ref{N5}. The decomposition of
the regular elliptic terms, the unipotent coefficients, the trace on the
theta functions and the stability of the principal series characters are
proved without it. It is used for the equivariance of the normalized
weighted distributions $I_A$ (Proposition~\ref{prop:IA-gl2}(a),~(b)), and
through it in the statement of part~(c) about $I_A$ and in the invariance
of the individual $I_{\mathcal O}$ and $I_\chi$ under $G_1(k)$. Parts~(a)
and~(b) as displayed do not use it.

\begin{theorem}\label{thm:stabler}
Let $S$, $f_S$ and $f=f_S\otimes e^S$ be as above. In part~(c)
assume~\ref{N5}.
\begin{enumerate}[label=(\alph*)]
\item (Geometric side.)
  \begin{multline*}
  I(f)=S_{\mathrm{ell}}(f)+S_{\mathrm{ell}}(f_-)
  +c_k\sum_{\{t,t^{-1}\}}I_A\bigl(a(t)_S,f_S\bigr)\\
  +\sum_{z=\pm1}\Bigl(\mathrm{vol}\bigl(G(\Q)\backslash\GA\bigr)f_S(z)+
  \frac{c_k}2\,I_A(z,f_S)+\frac12\sum_{\kappa\in\mathfrak X_S}
  \ell_S(\kappa)\,J^\kappa_{\mathrm{unip}}(z,f_S)\Bigr),
  \end{multline*}
  where $S_{\mathrm{ell}}$ is the stable distribution of
  Theorem~\ref{thm:elliptic-stable}, $t$ runs over $\OO_S^\times
  \smallsetminus\{\pm1\}$ modulo inversion, $\ell_S(1)=c_S$ and
  $\ell_S(\kappa)=L^S(1,\kappa)$ for $\kappa\ne1$.
\item (Spectral side.) $I(f)$ is the right side of
  Corollary~\ref{cor:itf-final}. The part of the discrete sum carried by
  the theta functions is $\frac12\sum_b\bigl(\Theta^{\mathrm{st}}_b(f)+
  \Theta^{\mathrm{st}}_b(f_-)\bigr)$, the sum over the classes $b\in k^\times/
  k^{\times2}$ with $\Theta_b$ pairwise orthogonal; the terms of the
  continuous spectrum are stable.
\item (Equivariance.) The individually indexed distributions occurring
  in~(a) transform under $G_1(k_S)$, acting by $f_S\mapsto f_S^h$, by the
  following characters of $\det h=a$: $\eta_\gamma$ for $J^{\mathrm{st}}_G(\gamma,\cdot)$ and
  $\eta_\gamma\kappa_E$ for $J^\kappa_G(\gamma,\cdot)$; $(\cdot,t)_S$ for
  $I_A(a(t)_S,\cdot)$; $(\cdot,z)_S$ for $f_S(z)$ and $I_A(z,\cdot)$;
  $\kappa_S\cdot(\cdot,z)_S$ for $J^\kappa_{\mathrm{unip}}(z,\cdot)$. The
  characters differ from one index to the next, so the sums displayed
  in~(a) are in general not eigendistributions; in particular
  $S_{\mathrm{ell}}$ is stable for the twisted action and is not an
  eigendistribution for the ordinary one. Both
  sides are invariant under $h\in G_1(k)\cap\prod_{v\notin S_k}
  Z(k_v)GL(2,\OO_v)$.
\end{enumerate}
\end{theorem}

\begin{proof}
(a) Theorem~\ref{thm:itf}(a)--(c), with Theorem~\ref{thm:elliptic-stable}
for the elliptic classes and Proposition~\ref{prop:unipotent}(c) for the
central ones; in~(c) of Theorem~\ref{thm:itf} the class of $a(t)$
contributes only if $t$ is a unit outside $S_k$. (b)
Proposition~\ref{prop:theta-packet} and the remark after it; the spaces
$\Theta_b$ for inequivalent $b$ are orthogonal because their irreducible
constituents are pairwise inequivalent (the constituents of
$\omega_{\psi'_v}$ and $\omega_{\psi''_v}$ are inequivalent when
$\psi''\ne\psi'_{c^2}$, by their restrictions to $N$). (c)
Lemma~\ref{lem:kappa-orbital}(b), Proposition~\ref{prop:IA-gl2}(b),
Lemma~\ref{lem:Ad}(d) and Proposition~\ref{prop:unipotent}(d). For the
last statement, such an $h$ acts on $f=f_S\otimes e^S$ through its
component in $G_1(k_S)$, by Lemma~\ref{lem:Ad}(e), and
Lemma~\ref{lem:spectrum-gl2}(a) applies.
\end{proof}

\paragraph{What the theorem does not say.} A stable trace formula for the
cover would be an identity
\[
I(f)=S(f)+S(f_-)
\]
with $S$ a stable distribution, and then an identification of $S(f)$
with $\frac12$ of the stable trace formula of $SO(3)$ evaluated at a
transfer of $f$; this is Li's theorem~\cite{Li21}. Parts~(a) and~(b)
exhibit this form for the regular elliptic terms and for the theta
functions. What is missing, in increasing order of difficulty, is:
\begin{enumerate}[label=(\roman*)]
\item the limit formulas that attach the terms $J^\kappa_{\mathrm{unip}}
  (z,\cdot)$, $\kappa\ne1$, to the singular elliptic points of the tori,
  and so distribute them between $S(f)$ and $S(f_-)$
  (Section~\ref{sec:unipotent});
\item the stabilization of the weighted terms $I_A(a(t),\cdot)$ and
  $I_A(z,\cdot)$, which are equivariant by
  Proposition~\ref{prop:IA-gl2} but not stable;
\item the transfer $f\mapsto f'$ to $SO(3)$ and the corresponding
  fundamental lemma~\cite{Li11,Li12};
\item on the spectral side, the description of the cuspidal spectrum by
  packets~\cite{Wa80,Wa91,GI}.
\end{enumerate}
The equivariance in~(c) is the property that the 1995 notes called
``$\kappa$-stable''. It is weaker than stability, and the part of $I(f)$
that is invariant under all of $G_1(k_S)$, which the 1995 notes took as
the stable part, is small: by Lemma~\ref{lem:kappa-orbital}(b) and
Proposition~\ref{prop:IA-gl2}(b), a regular elliptic or hyperbolic
$\gamma$ contributes to it only if, at each place of $S_k$, $D^+(\gamma)$
is a square or $\theta$ times a square; for instance if $\gamma$ is a
square in its torus.

\subsection{Status of the results}\label{sec:status7}

\begin{remarkrev}\label{rem:sec7-status}
\emph{What rests on citation.} Proved here, given the earlier sections:
Lemmas~\ref{lem:Ad} and~\ref{lem:eta} (with the cocycle~\eqref{eq:sigma-gl2}
on $GL(2)$ from~\cite{K2}; part~(c) of Lemma~\ref{lem:Ad} is also checked
directly); Lemma~\ref{lem:Delta}; Lemma~\ref{lem:weil-kernel};
Lemma~\ref{lem:kappa-orbital}; Theorem~\ref{thm:elliptic-stable}, with
the Hasse norm theorem and the index of the global norm group
from~\cite{CF}; Lemma~\ref{lem:variance-gl2};
Proposition~\ref{prop:IA-gl2}, parts~(a) and~(b) under~\ref{N5};
Proposition~\ref{prop:unipotent}; Proposition~\ref{prop:theta-packet}(b),
(c) given~(a); and Proposition~\ref{prop:weil-char} on the regular set,
with Lemma~\shellref. Taken from the literature: at the
archimedean places, the standard fact that an operator with a Schwartz
kernel is of trace class with trace the integral over the diagonal; and
at all places, that the character of the Weil representation is a locally
integrable function on the whole group~\cite{Mak,Th08}, which is used to call $\Theta^+_\psi-\Theta^-_\psi$ a
stable character and in the second formula for $\Theta^{\mathrm{st}}_b$ in
Section~\ref{sec:spectral-gl2}. Formula~\eqref{eq:weil-char-full} agrees
with~\cite[Theorem~1C]{Th08} up to the sign left open there, so its sign
rests on our computation alone; we have not compared it with the formulas
of~\cite{Mak,Th13} in detail. Also from the
literature: the convergence and square-integrability of theta functions
in Proposition~\ref{prop:theta-packet}(a), Tate's local zeta
integrals~\cite{Ta}, and the statements about Waldspurger's packets after
Proposition~\ref{prop:theta-packet}, which are not used.

\emph{What is not done.} The condition~\ref{N5} is a hypothesis at the
places where the normalizing factors are not explicit: we have not shown
that factors satisfying it exist there, nor that those of~\cite{LiII} do
(Remark~\ref{rem:N5}). The statements in the text about the relation
with the endoscopy of~\cite{Li11,Li15,Li21} are comparisons of shape: we
have not identified $J^{\mathrm{st}}_G$ with Li's stable orbital integral
beyond Proposition~\ref{prop:weil-char}, nor $\Delta_\psi$ with his
transfer factor, nor checked the normalizations in
$\iota(\tilde G,H)=\frac12$. The limit formulas of
Section~\ref{sec:unipotent} at the singular elliptic points, the
stabilization of the weighted terms, the transfer to $SO(3)$ and the
description of the cuspidal spectrum are not treated. Sections~\ref{sec:background}--\ref{sec:invariant}
and this section were refereed independently (see the Acknowledgements). An unpublished
stabilization of the trace formula for covers of $SL(2)$ by Hiraga and
Ikeda is mentioned in~\cite{Li15}; see also~\cite{HI13}.
\end{remarkrev}

\begin{remarkrev}\label{rem:verified-sec7}
The script \texttt{checks/sec7\_n2\_check.py} verifies, in exact
arithmetic at $v=\infty,2,3,5,7$ on random elements of $SL(2,\Q)$ and
$GL(2,\Q)$: Lemma~\ref{lem:Ad}(c),(d), both as a homomorphism for
$\beta_v$ and as conjugation in the cover of $GL(2)$;
Lemma~\ref{lem:eta} on the tori $T_\theta(\Q)$ and on the split torus,
and Example~\ref{ex:eta}(d); the conjugation invariance of $\Delta_\psi$,
Lemma~\ref{lem:Delta}(b),(c), and Example~\ref{ex:Delta-real}; the
identity between Weil indices behind~\eqref{eq:weil-kernel}; the
consistency of~\eqref{eq:weil-char} with~\eqref{eq:weil-char-full} under
translation by $(-1,1)$, and the agreement of~\eqref{eq:weil-char-full}
with~\cite[Theorem~1C]{Th08} up to sign;
identity~\eqref{eq:eta-minus}; and Lemma~\ref{lem:variance-gl2}(a).
Numerically: Example~\ref{ex:Delta-finite} for $q=3,5,7,11$;
Proposition~\ref{prop:unipotent}(a) for $k=\Q$, $S=\{\infty,2\}$ at
$s=3$, the local identity~\eqref{eq:squares} at $2$ and $3$, and the
constants of Example~\ref{ex:unip-coeff}, including the invariance of the
coefficients under $\OO_S^\times$; Lemma~\shellref at $2$, $3$ and $5$; the two facts used in the proof of
Proposition~\ref{prop:theta-packet}(b); and the parity in $\theta$ of the
sum and the difference of the two discrete series characters in the
paragraph after that proposition.
\end{remarkrev}
